\documentclass[11pt,a4paper,openany,twoside=false]{book}

% ============================================================
% Paquetes básicos
% ============================================================
\usepackage[utf8]{inputenc}
\usepackage[T1]{fontenc}
\usepackage[spanish,es-tabla]{babel}
\usepackage{lmodern}
\usepackage[a4paper,margin=2.5cm,headheight=14pt]{geometry}

% ============================================================
% Gráficos y figuras
% ============================================================
\usepackage{graphicx}
\usepackage{float}
\usepackage{caption}
\usepackage{subcaption}
\usepackage{tikz}
\usetikzlibrary{positioning,arrows.meta,shapes.geometric,calc,fit,backgrounds,decorations.pathreplacing,shapes.misc}
\usepackage{tikzscale}

% ============================================================
% Tablas
% ============================================================
\usepackage{booktabs}
\usepackage{longtable}
\usepackage{array}
\usepackage{multirow}
\usepackage{tabularx}
\usepackage{makecell}

% ============================================================
% Matemática
% ============================================================
\usepackage{amsmath}
\usepackage{amssymb}
\usepackage{amsfonts}
\usepackage{mathtools}
\usepackage{bm}

% ============================================================
% Código (minted + listings)
% ============================================================
\usepackage{minted}
\usepackage{listings}
\setminted{
    fontsize=\footnotesize,
    linenos,
    numbersep=8pt,
    frame=lines,
    framesep=2mm,
    tabsize=4,
    baselinestretch=1.05,
    breaklines=false,
    breakbytoken=false,
    python3=true,
    style=manni,
    fontfamily=tt,
    fontshape=tt,
    fontseries=bf,
    obeytabs=true,
    showtabs=false,
    tab=\rightarrowfill
}
% Listings como fallback (no necesita shell-escape)
\lstset{
    basicstyle=\ttfamily\footnotesize,
    keywordstyle=\color{blue!70!black}\bfseries,
    commentstyle=\color{green!50!black},
    stringstyle=\color{red!70!black},
    numberstyle=\tiny\color{gray},
    numbers=left,
    numbersep=8pt,
    frame=single,
    framesep=4pt,
    tabsize=4,
    breaklines=false,
    showstringspaces=false,
    extendedchars=true,
    inputencoding=utf8,
    literate=
        {á}{{\'a}}1 {é}{{\'e}}1 {í}{{\'i}}1 {ó}{{\'o}}1 {ú}{{\'u}}1
        {Á}{{\'A}}1 {É}{{\'E}}1 {Í}{{\'I}}1 {Ó}{{\'O}}1 {Ú}{{\'U}}1
        {ñ}{{\~n}}1 {Ñ}{{\~N}}1 {¿}{{?`}}1 {¡}{{!`}}1
        {ü}{{\"u}}1 {Ü}{{\"U}}1
        {α}{{$\alpha$}}1 {β}{{$\beta$}}1 {γ}{{$\gamma$}}1
        {→}{{$\rightarrow$}}1 {←}{{$\leftarrow$}}1,
    language=Python,
    aboveskip=8pt,
    belowskip=8pt
}

% ============================================================
% Hipervínculos
% ============================================================
\usepackage[colorlinks=true,linkcolor=blue!60!black,urlcolor=blue!60!black,citecolor=blue!60!black,bookmarksopen=true,bookmarksnumbered=true]{hyperref}
\usepackage{cleveref}
\crefname{section}{Sección}{Secciones}
\crefname{chapter}{Capítulo}{Capítulos}
\crefname{figure}{Figura}{Figuras}
\crefname{table}{Tabla}{Tablas}
\crefname{listing}{Listado}{Listados}
\crefname{equation}{Ecuación}{Ecuaciones}

% ============================================================
% Encabezados
% ============================================================
\usepackage{fancyhdr}
\pagestyle{fancy}
\fancyhf{}
\fancyhead[L]{\small\textit{TextReIDNet}}
\fancyhead[R]{\small\textit{\nouppercase{\leftmark}}}
\fancyfoot[C]{\thepage}
\renewcommand{\headrulewidth}{0.4pt}

% ============================================================
% Utilidades
% ============================================================
\usepackage{xcolor}
\definecolor{codekw}{RGB}{0,0,180}
\definecolor{codestr}{RGB}{163,21,21}
\definecolor{codebg}{RGB}{248,248,248}
\usepackage{enumitem}
\setlist{nosep,leftmargin=*}
\usepackage{csquotes}
\usepackage{titlesec}
\titleformat{\chapter}[display]
    {\normalfont\huge\bfseries}
    {\chaptertitlename~\thechapter}{20pt}{\Huge}
\titleformat{\section}
    {\normalfont\Large\bfseries}
    {\thesection}{1em}{}
\titleformat{\subsection}
    {\normalfont\large\bfseries}
    {\thesubsection}{1em}{}

% ============================================================
% Información del documento
% ============================================================
\title{\Huge\textbf{TextReIDNet}\\[0.3cm]
\Large Diseño e implementación de un modelo\\[0.2cm]
de reidentificación de personas\\[0.2cm]
basado en texto\\[1cm]
\normalfont\large Documento técnico — Proyecto \texttt{textreid-train}}
\author{}
\date{\today}

% ============================================================
% Comandos personalizados
% ============================================================
\newcommand{\todo}[1]{\textcolor{red}{\textbf{[TODO: #1]}}}
\newcommand{\dims}[1]{\textcolor{blue!70!black}{\texttt{(#1)}}}
\newcommand{\param}[2]{\texttt{\textbf{#1}}\index{#1} \hfill \textit{#2}}
\newcommand{\file}[1]{\texttt{\detokenize{#1}}}
\DeclareRobustCommand{\fline}[2]{\detokenize{#1}:\textbf{#2}}
\newcommand{\reid}{re-ID}
\newcommand{\gru}{GRU}
\newcommand{\dsc}{DSC}
\newcommand{\eff}{EfficientNet-B0}
\newcommand{\shapesym}[1]{\textcolor{blue!70!black}{\texttt{#1}}}
\DeclareRobustCommand{\acronym}[2]{\textsc{#1}\index{#1}}

\begin{document}

% ============================================================
% Portada y frontmatter
% ============================================================
\frontmatter
\maketitle

\begin{center}
\textbf{\large Resumen}
\end{center}

\noindent
Este libro documenta en profundidad el proyecto \texttt{textreid-train}, una
implementación limpia y enfocada al entrenamiento del modelo \textbf{TextReIDNet}
para \emph{reidentificación de personas basada en texto}. El modelo combina
un backbone visual \eff{} con una red de lenguaje basada en \emph{BERT} +
\gru{} bidireccional, fusionando ambas modalidades en un espacio de
\emph{embedding} conjunto de dimensión 1024. El entrenamiento se realiza
sobre el dataset \textbf{CUHK-PEDES}.

El documento parte del paper original de Agyeman \emph{et al.} (\emph{IEEE
Access}, 2024) -- ``Reidentificación descentralizada de personas basada
en texto en redes multicámara'' -- cuya versión traducida al español se
encuentra en \file{docs/paper.md}, y desciende hasta la última línea de
cada archivo del repositorio: explica qué hace cada clase, qué forma
tienen los tensores en cada capa, dónde se define cada parámetro de
configuración y por qué se eligieron esos valores.

Está pensado como material de referencia para quien quiera entender,
modificar, depurar o reentrenar el modelo. No se asume experiencia
previa con \emph{re-ID}, pero sí familiaridad básica con \emph{PyTorch}
y \emph{deep learning}.

\vspace{1cm}

\tableofcontents
\listoffigures
\listoftables
\listoflistings

% ============================================================
% Capítulos
% ============================================================
\mainmatter

% ============================================================
% Capítulo 0 — Prólogo
% ============================================================
\chapter{Prólogo}

Este libro es una guía exhaustiva del proyecto \texttt{textreid-train}:
un repositorio limpio, enfocado exclusivamente al entrenamiento y la
evaluación del modelo \textbf{TextReIDNet} para \emph{reidentificación
de personas basada en texto}. No es un manual de uso ni un
\emph{quickstart}; es un documento de referencia técnica que recorre
cada archivo del proyecto y explica, capa por capa, qué hace y por qué.

\section{¿A quién está dirigido?}

\begin{itemize}
  \item \textbf{Estudiantes o practitioners de visión por computador}
        que quieran entender un sistema real de
        \emph{re-ID} basado en texto, desde la matemática del paper
        hasta los detalles de implementación.
  \item \textbf{Investigadores} que busquen reproducir, comparar o
        extender el modelo.
  \item \textbf{Ingenieros} que tengan que mantener, depurar o
        desplegar este código en producción.
\end{itemize}

\section{Qué asume el lector}

\begin{itemize}
  \item Familiaridad básica con \textbf{PyTorch} (tensores,
        \texttt{nn.Module}, \texttt{DataLoader}).
  \item Conocimiento general de \emph{deep learning}: redes
        convolucionales, recurrentes, funciones de pérdida y
        optimización.
  \item Conocimiento elemental de procesamiento de lenguaje natural:
        \emph{tokenización}, \emph{embeddings}, \acronym{GRU}{GRU}.
\end{itemize}

No se asume experiencia previa con \emph{reidentificación de personas},
\emph{EfficientNet}, \emph{BERT} ni métricas como \acronym{mAP}{mAP}.

\section{Convenciones tipográficas}

\begin{itemize}
  \item Los nombres de archivos se escriben en
        \file{\texttt{monoespaciado}} y se citan con su ruta relativa
        al raíz del proyecto: \file{model/textreidnet.py}.
  \item Las referencias a líneas específicas se hacen con el formato
        \fline{archivo.py}{L}. Por ejemplo, \fline{config.py}{42}
        apunta a la línea 42 de \file{config.py}.
  \item Las formas de los tensores se muestran entre paréntesis en
        color azul: \dims{(B, 1280, 16, 16)} significa
        \emph{batch} × \emph{canales} × \emph{alto} × \emph{ancho}.
  \item Los parámetros de configuración aparecen en
        \texttt{\textbf{negrita mono}}.
  \item Los listados de código van numerados y muestran sólo la región
        relevante. La numeración coincide con la del archivo
        original.
\end{itemize}

\section{Estructura del libro}

El libro se divide en cuatro partes:

\begin{enumerate}
  \item \textbf{Fundamentos teóricos} (Caps.~\ref{cap:introduccion}
        y~\ref{cap:arquitectura-teorica}): contexto del problema,
        formulación matemática y arquitectura de U-TextReIDNet según
        el paper.
  \item \textbf{Sistema descentralizado} (Cap.~\ref{cap:sistema-descentralizado}):
        cómo se despliega el modelo en una red de cámaras reales.
  \item \textbf{Implementación} (Caps.~\ref{cap:dataset}--\ref{cap:evaluacion}):
        dataset, configuración, redes, pérdidas, entrenamiento y
        evaluación. Es el grueso del libro.
  \item \textbf{Apéndices} (Cap.~\ref{cap:apendices}): glosario,
        tablas de referencia rápida.
\end{enumerate}

\section{Código fuente}

Todos los listados del libro provienen de los archivos reales del
repositorio \texttt{textreid-train}. Se recomienda abrir el repositorio
en paralelo a la lectura para poder navegar entre el PDF y el código
sin perderse.

\begin{center}
\begin{tikzpicture}
  \node[draw, rounded corners, fill=codebg, minimum width=6cm, align=center]
    {\textbf{\large textreid-train/}\\[2pt]
     \texttt{model/} · \texttt{datasets/} · \texttt{evaluation/}\\[1pt]
     \texttt{utils/} · \texttt{scripts/} · \texttt{config.py}};
\end{tikzpicture}
\end{center}

\section*{Sobre este documento}

\noindent\textbf{Versión}: 1.0\\
\textbf{Fecha de compilación}: \today\\
\textbf{Paper de referencia}: Agyeman \emph{et al.}, ``Decentralized
Text-Based Person Re-Identification in Multicamera Networks'',
\emph{IEEE Access}, noviembre de 2024. DOI:~\href{https://doi.org/10.1109/ACCESS.2024.3501382}{10.1109/ACCESS.2024.3501382}.\\
\textbf{Traducción del paper al español}:
\file{docs/paper.md} (552 líneas).

\vspace{1em}
\noindent\rule{\textwidth}{0.4pt}
\vspace{0.5em}

\noindent\textit{Conviene leer este libro con un navegador de código
abierto. Cada vez que veas una referencia como
\fline{model/textreidnet.py}{13}, podrás saltar directamente al
constructor de la clase \texttt{TextReIDNet}.}

% ============================================================
% Capítulo 1 — Introducción
% ============================================================
\chapter{Introducción a la reidentificación de personas}
\label{cap:introduccion}

Este capítulo sienta las bases del problema que el modelo intenta
resolver. Toda la información proviene del paper traducido al español
en \file{docs/paper.md}.

\section{¿Qué es la reidentificación de personas?}

La \emph{reidentificación de personas} (en inglés \emph{person
re-identification}, abreviada \reid{}) es la tarea de reconocer a la
\emph{misma} persona en distintas fotografías capturadas por cámaras
diferentes -- o por la misma cámara en distintos momentos. La consulta
de búsqueda puede ser:

\begin{itemize}
  \item una \textbf{imagen} (\emph{image-to-image re-ID}),
  \item una \textbf{secuencia de vídeo},
  \item una \textbf{descripción textual} (\emph{text-based re-ID}).
\end{itemize}

A diferencia del reconocimiento facial, la \reid{} no depende de
características faciales -- que pueden ser de baja calidad o estar
restringidas por privacidad -- sino de atributos visibles como ropa,
color, tipo de cuerpo o accesorios \cite{agyeman2024decentralized}.
Esto la hace especialmente útil en entornos concurridos donde es
difícil obtener caras nítidas.

\section{Reidentificación basada en texto}

La \reid{} basada en texto (\emph{text-based person re-identification})
consiste en recuperar la imagen de una persona a partir de una
descripción en lenguaje natural. Por ejemplo, ante la consulta:

\begin{quote}
\textit{``Una mujer lleva una camiseta negra con letras rosas. Viste
pantalones cortos color caqui y zapatos blancos.''}
\end{quote}

\noindent el sistema debe devolver, de una galería de miles de
imágenes, las que correspondan a esa persona.

Esta modalidad tiene dos ventajas prácticas frente a la \reid{}
basada en imagen:

\begin{enumerate}
  \item Las consultas son más \textbf{accesibles}: cualquiera puede
        describir a una persona, pero no siempre se tiene una foto
        reciente.
  \item Resuelve \textbf{problemas de privacidad}: en muchas
        jurisdicciones capturar y almacenar caras está restringido,
        mientras que las descripciones textuales son menos invasivas.
\end{enumerate}

\section{Desafíos de los sistemas existentes}

Los algoritmos \emph{state-of-the-art} de \reid{} por texto suelen
ser modelos enormes (cientos de millones de parámetros) que
priorizan precisión sobre eficiencia. Esto trae tres problemas en
aplicaciones reales:

\begin{itemize}
  \item \textbf{Requisitos computacionales altos}: no se pueden
        ejecutar en dispositivos embebidos como una cámara
        inteligente.
  \item \textbf{Dependencia de arquitecturas centralizadas}: todas
        las cámaras envían vídeo a un servidor central que ejecuta
        la identificación. Esto genera cuellos de botella,
        problemas de escalabilidad y consumo elevado de ancho de
        banda.
  \item \textbf{Imágenes pre-recortadas}: la mayoría de los
        algoritmos asumen que la imagen de entrada contiene una
        sola persona ya recortada, lo cual es irreal en escenas
        concurridas.
\end{itemize}

\section{La solución propuesta: U-TextReIDNet}

El paper de Agyeman \emph{et al.}~\cite{agyeman2024decentralized}
propone \textbf{U-TextReIDNet} (Unified Text Re-Identification
Network), un modelo:

\begin{itemize}
  \item \textbf{Preciso}: alcanza \textbf{54,02\%} Top-1 en
        CUHK-PEDES y \textbf{38,45\%} Top-1 en RSTPReid.
  \item \textbf{Ligero}: solo \textbf{38,77 millones} de
        parámetros totales.
  \item \textbf{Eficiente}: corre en tiempo real en un
        \emph{NVIDIA Jetson Nano}.
  \item \textbf{Unificado}: integra la detección de personas con
        la identificación, eliminando la necesidad de pre-recortar.
\end{itemize}

A diferencia de sus predecesores \emph{MHParsNet} (detección) y
\emph{TextReIDNet} (identificación), que eran módulos separados,
U-TextReIDNet los une en un solo modelo end-to-end.

\section{El dataset CUHK-PEDES}

Este proyecto entrena sobre \textbf{CUHK-PEDES}
\cite{li2017personsearch}, uno de los benchmarks estándar para
\reid{} por texto. Estadísticas según el paper:

\begin{table}[H]
\centering
\caption{Estadísticas del dataset CUHK-PEDES original.}
\label{tab:cuhkpedes-stats}
\begin{tabular}{lrrr}
\toprule
\textbf{Split} & \textbf{Imágenes} & \textbf{Descripciones} & \textbf{Personas} \\
\midrule
Train & 34\,054 & 68\,126 & 11\,003 \\
Val   &  3\,078 &  6\,158 &  1\,000 \\
Test  &  3\,074 &  6\,156 &  1\,000 \\
\midrule
\textbf{Total} & \textbf{40\,206} & \textbf{80\,440} & \textbf{13\,003} \\
\bottomrule
\end{tabular}
\end{table}

\noindent Cada imagen tiene asociadas \emph{dos} descripciones
textuales escritas por anotadores humanos. En este proyecto se usa
una versión alternativa descargada de \emph{HuggingFace}
(\texttt{PeterPanTheGenius/CUHK-PEDES}) que tiene una sola
descripción por imagen pero está disponible sin necesidad de
solicitar acceso por email. Los detalles están en el
\cref{cap:dataset}.

\section{Lo que este libro no cubre}

U-TextReIDNet completo (paper) incluye dos componentes:

\begin{itemize}
  \item \textbf{R-MHParsNet}: detector de múltiples personas
        (\emph{FPN} + EfficientNet-B0 + cabezas de segmentación).
        Tiene \textbf{6,49M} parámetros y NO está implementado en
        este proyecto.
  \item \textbf{TextReIDNet}: el modelo de identificación
        imagen-texto (EfficientNet-B0 + BERT + BiGRU). Éste es el
        núcleo del proyecto \texttt{textreid-train}.
\end{itemize}

Este libro se centra exclusivamente en \textbf{TextReIDNet} y su
entrenamiento. Para la parte descentralizada (Jetson Nano, ZMQ,
MQTT, Flask) se da contexto en el \cref{cap:sistema-descentralizado}
pero sin entrar en código, ya que el repositorio no la incluye.

% ============================================================
% Capítulo 2 — Arquitectura teórica U-TextReIDNet
% ============================================================
\chapter{U-TextReIDNet: arquitectura teórica}
\label{cap:arquitectura-teorica}

Este capítulo desglosa la arquitectura completa del modelo
U-TextReIDNet tal como se describe en \file{docs/paper.md}
(§\textsc{Iii}-A, líneas 146--178). Aunque el repositorio
\texttt{textreid-train} implementa solo la parte de identificación
(\emph{TextReIDNet}), entender la versión unificada ayuda a
comprender las elecciones de diseño.

\section{Formulación del problema}

Dado un fotograma de vídeo con $n$ personas distintas
$\mathbf{V} \in \mathbb{R}^{H \times W \times 3}$ y una descripción
textual $\mathbf{T}$ en lenguaje natural, los objetivos son:

\begin{enumerate}
  \item \textbf{Detectar y segmentar} las $n$ instancias humanas en
        $\mathbf{V}$, produciendo una galería
        $\mathbf{I} = \{\mathbf{I}_1, \ldots, \mathbf{I}_n\}$.
        Si $n = 0$, la galería está vacía.
  \item \textbf{Identificar} la imagen $\mathbf{I}^* \in \mathbf{I}$
        que mejor se ajusta a $\mathbf{T}$ mediante una función de
        coincidencia multimodal $\mathcal{M}$.
\end{enumerate}

La función $\mathcal{M}$ usa las \emph{incrustaciones} (embeddings)
visuales y textuales:

\begin{align}
  \mathbf{I}' &= F_i(\mathbf{I}) && \text{embedding visual} \\
  \mathbf{T}' &= F_t(\mathbf{T}) && \text{embedding textual}
\end{align}

\noindent y calcula su similitud con la métrica del coseno:

\begin{equation}
  \text{sim}(\mathbf{I}', \mathbf{T}') =
  \frac{\mathbf{I}' \cdot \mathbf{T}'}{\|\mathbf{I}'\|\,\|\mathbf{T}'\|}
  \label{eq:cosine}
\end{equation}

\noindent El resultado óptimo es la imagen que maximiza esa
similitud.

\section{Componentes principales}

U-TextReIDNet se compone de tres bloques:

\begin{enumerate}
  \item \textbf{Red visual}: \emph{backbone} EfficientNet-B0 +
        capa de convolución separable en profundidad.
  \item \textbf{Red de lenguaje}: BERT + BiGRU bidireccional.
  \item \textbf{Componente de coincidencia}: espacio de embedding
        conjunto + pérdidas.
\end{enumerate}

\subsection{Red visual}

\paragraph{Detector R-MHParsNet (no implementado en este proyecto).}

R-MHParsNet es una FPN (Feature Pyramid Network) basada en
EfficientNet-B0, con cabezas de segmentación para generar máscaras
de instancia humana. Genera los bounding boxes
$\{(x_1,y_1,x_2,y_2)\}_{i=1}^{n}$ que delimitan cada persona
detectada. Tiene 6,49M parámetros (un 79,9\% más pequeño que su
predecesor MHParsNet).

\paragraph{Generador de embedding visual (sí implementado).}

Para cada persona detectada, recortada como
$\mathbf{I} \in \mathbb{R}^{384 \times 128 \times 3}$:

\begin{enumerate}
  \item Se redimensiona a $\mathbb{R}^{3 \times 384 \times 128}$
        (HxW o WxH según convención; el paper dice 384 alto × 128
        ancho).
  \item Pasa por EfficientNet-B0 y se extrae la salida de la
        \textbf{etapa 9} (la última): un tensor de
        $\mathbb{R}^{1280 \times 12 \times 4}$. Los 1280 canales
        son la profundidad de features; $12 \times 4$ son las
        dimensiones espaciales (es decir, $384 / 32 = 12$ y
        $128 / 32 = 4$).
  \item Una capa de convolución separable en profundidad
        (\dsc{}) con 1024 filtros $3 \times 3$ reduce a
        $\mathbb{R}^{1024 \times 12 \times 4}$.
\end{enumerate}

\noindent\textbf{Nota técnica:} en este proyecto las imágenes
entran a $512 \times 512$ (ver \fline{datasets/bases.py}{43} y
\cref{cap:dataset}), no a $384 \times 128$. Eso cambia las
dimensiones espaciales en el resto del pipeline. Se conserva la
misma profundidad de canales.

\subsection{Red de lenguaje}

El componente de lenguaje procesa la descripción textual
$\mathbf{T}$:

\begin{enumerate}
  \item \textbf{BERT} tokeniza y recorta/padea a una longitud
        fija de 100 tokens: $(z_{+1}, \ldots, z_{+100})$.
        El tokenizer es \texttt{bert-base-uncased}.
  \item Una \textbf{capa de embedding} de 512 dimensiones convierte
        cada token en un vector denso:
        $(\mathbf{e}_1, \ldots, \mathbf{e}_{100}) \in
        \mathbb{R}^{100 \times 512}$.
  \item Una \textbf{GRU bidireccional} con 1024 unidades ocultas
        procesa la secuencia y produce:
        \begin{equation}
          \mathbf{T}' = (\overrightarrow{\text{GRU}} +
                          \overleftarrow{\text{GRU}}) / 2
          \in \mathbb{R}^{1024 \times 100 \times 1}
          \label{eq:bigru}
        \end{equation}
        La suma promedio de las direcciones hacia adelante y hacia
        atrás captura el contexto en ambos sentidos.
\end{enumerate}

\subsection{Componente de coincidencia de embedding conjunto}

Las dos modalidades se proyectan a un \textbf{espacio conjunto} de
dimensión 1024:

\begin{itemize}
  \item $\mathbf{I}' \in \mathbb{R}^{1024 \times 12 \times 4}$
        (visual).
  \item $\mathbf{T}' \in \mathbb{R}^{1024 \times 100 \times 1}$
        (textual).
\end{itemize}

Para optimizar el \emph{matching} se usan dos pérdidas:

\begin{itemize}
  \item \textbf{Pérdida de identidad} ($\mathcal{L}_{\text{ID}}$):
        cross-entropy que refuerza que cada persona tenga una
        representación única en el espacio.
  \item \textbf{Pérdida de ranking/clasificación}
        ($\mathcal{L}_{\text{rank}}$): basada en triplet loss,
        garantiza que las coincidencias correctas obtengan
        puntuaciones más altas que las incorrectas.
\end{itemize}

\noindent La suma ponderada es la pérdida total:
\begin{equation}
  \mathcal{L} = \alpha \cdot \mathcal{L}_{\text{rank}}
              + \beta  \cdot \mathcal{L}_{\text{ID}}
\end{equation}

\noindent En este proyecto $\alpha = \beta = 1.0$
(\fline{config.py}{87-88}).

\section{Diagrama de bloques}

\begin{figure}[H]
\centering
% ============================================================
% Diagrama TikZ: Arquitectura U-TextReIDNet
% ============================================================
\begin{tikzpicture}[
    font=\sffamily\small,
    node distance=0.6cm and 0.5cm,
    every node/.style={align=center},
    block/.style={
        rectangle, draw=black, thick, rounded corners=2pt,
        minimum height=0.9cm, minimum width=2cm,
        fill=blue!8
    },
    block2/.style={
        rectangle, draw=black, thick, rounded corners=2pt,
        minimum height=0.9cm, minimum width=2cm,
        fill=orange!15
    },
    block3/.style={
        rectangle, draw=black, thick, rounded corners=2pt,
        minimum height=0.9cm, minimum width=2cm,
        fill=green!15
    },
    impl/.style={
        rectangle, draw=blue!60!black, very thick, rounded corners=2pt,
        minimum height=0.9cm, minimum width=2cm,
        fill=blue!15
    },
    arr/.style={
        -{Stealth[length=2.5mm,width=2mm]}, thick
    },
    group/.style={
        rectangle, draw=gray, dashed, thick, rounded corners=4pt,
        inner sep=10pt, fill=gray!5
    }
]

% ======================== RAMA VISUAL ========================
\node[block2] (v_in)  {Frame\\$V \in \mathbb{R}^{H \times W \times 3}$};
\node[block, below=of v_in] (rmhp)  {R-MHParsNet\\(detector)};
\node[block, below=of rmhp] (crop)  {Crop + resize\\$384 \times 128$};
\node[impl, below=of crop] (eff)   {\textbf{EffNet-B0}\\stage 9};
\node[impl, below=of eff]  (dsc1)  {\textbf{DSC}\\1280 $\to$ 1024};
\node[impl, below=of dsc1] (amp)   {Adaptive\\MaxPool 1$\times$1};
\node[impl, below=of amp]  (dsc2)  {\textbf{DSC}\\1024 $\to$ 1024};
\node[impl, below=of dsc2] (ve)    {$\mathbf{I}' \in \mathbb{R}^{1024}$};

% ======================== RAMA TEXTO ========================
\node[block2, right=4cm of v_in] (t_in)  {Texto\\$T$};
\node[impl, right=4cm of rmhp] (bert) {\textbf{BERT}\\100 tokens};
\node[impl, right=4cm of crop] (emb)  {\textbf{Embedding}\\$\mathbb{R}^{512}$};
\node[impl, right=4cm of eff]  (gru)  {\textbf{BiGRU}\\1024 ocultas};
\node[impl, right=4cm of dsc1] (sum)  {Suma fwd+bwd /2};
\node[impl, right=4cm of amp]  (perm) {Permute\\+unsqueeze};
\node[impl, right=4cm of dsc2] (tmax) {Max sobre\\dimensión 2};
\node[impl, right=4cm of ve]   (dsct) {\textbf{DSC}\\1024 $\to$ 1024};
\node[impl, right=4cm of ve, yshift=-0.8cm] (te)  {$\mathbf{T}' \in \mathbb{R}^{1024}$};

% ======================== MATCHING ========================
\node[block3, below=2cm of ve] (cos)  {Similitud coseno\\$\mathcal{M}$};
\node[block3, right=4cm of cos] (rank) {$\mathcal{L}_{\text{rank}}$};
\node[block3, right=4cm of rank] (id)  {$\mathcal{L}_{\text{ID}}$};

% Conexiones rama visual
\draw[arr] (v_in) -- (rmhp);
\draw[arr] (rmhp) -- (crop);
\draw[arr] (crop) -- (eff);
\draw[arr] (eff)  -- (dsc1);
\draw[arr] (dsc1) -- (amp);
\draw[arr] (amp)  -- (dsc2);
\draw[arr] (dsc2) -- (ve);

% Conexiones rama texto
\draw[arr] (t_in) -- (bert);
\draw[arr] (bert) -- (emb);
\draw[arr] (emb)  -- (gru);
\draw[arr] (gru)  -- (sum);
\draw[arr] (sum)  -- (perm);
\draw[arr] (perm) -- (tmax);
\draw[arr] (tmax) -- (dsct);
\draw[arr] (dsct) -- (te);

% Matching
\draw[arr] (ve) -- (cos);
\draw[arr] (te) -- (cos);
\draw[arr] (cos.east) -| (rank.west);
\draw[arr] (cos.east) -| ([yshift=0.6cm]id.west) -- (id.west);
\draw[arr] (rank.east) -- (id.west);

% Agrupaciones
\begin{scope}[on background layer]
  \node[group, fit=(rmhp), label={[gray]above:R-MHParsNet (paper)}] {};
  \node[group, fit=(eff) (dsc1) (amp) (dsc2) (ve), label={[blue!60!black]above:\textbf{Implementado en este proyecto}}] {};
  \node[group, fit=(bert) (emb) (gru) (sum) (perm) (tmax) (dsct) (te), label={[blue!60!black]above:\textbf{Implementado en este proyecto}}] {};
\end{scope}

% Leyenda
\node[draw, rounded corners, fill=orange!15, minimum width=2cm, minimum height=0.5cm, anchor=north west] at (-9, -10) {};
\node[anchor=west] at (-7, -10) {Paper, no en código};
\node[draw, rounded corners, fill=blue!15, minimum width=2cm, minimum height=0.5cm, anchor=north west] at (-9, -10.8) {};
\node[anchor=west] at (-7, -10.8) {Implementado aquí};
\node[draw, rounded corners, fill=green!15, minimum width=2cm, minimum height=0.5cm, anchor=north west] at (-9, -11.6) {};
\node[anchor=west] at (-7, -11.6) {Funciones de pérdida};

\end{tikzpicture}

\caption{Arquitectura de U-TextReIDNet. Las cajas con relleno gris
son las que este proyecto implementa; las cajas con relleno claro
son las del paper que no están en el código (R-MHParsNet y sistema
distribuido).}
\label{fig:arquitectura}
\end{figure}

\section{Dimensiones en cada etapa: tabla resumen}

\begin{table}[H]
\centering
\caption{Dimensiones de los tensores en cada etapa del modelo.
\textit{B} denota el tamaño de batch. Las columnas ``Etapa paper''
y ``Etapa proyecto'' difieren únicamente en la resolución espacial
de entrada.}
\label{tab:dimensiones}
\small
\begin{tabular}{lrll}
\toprule
\textbf{Componente} & \textbf{Forma} & \textbf{Origen} & \textbf{Ubicación} \\
\midrule
Imagen de entrada
  & \dims{(B, 3, 512, 512)}
  & resize/colocar en canvas
  & \fline{datasets/bases.py}{128} \\
\midrule
EfficientNet-B0 stage 9
  & \dims{(B, 1280, 16, 16)}
  & salida \texttt{model.features}
  & \fline{model/visual_network.py}{29} \\
\midrule
DSC 1280 $\to$ 1024
  & \dims{(B, 1024, 16, 16)}
  & conv. separable en profundidad
  & \fline{model/textreidnet.py}{41} \\
\midrule
AdaptiveMaxPool2d(1,1)
  & \dims{(B, 1024, 1, 1)}
  & pool global
  & \fline{model/textreidnet.py}{42} \\
\midrule
DSC 1024 $\to$ 1024
  & \dims{(B, 1024, 1, 1)}
  & conv. separable en profundidad
  & \fline{model/textreidnet.py}{43} \\
\midrule
Visual embedding final
  & \dims{(B, 1024, 1)}
  & \texttt{squeeze(-1)}
  & \fline{model/textreidnet.py}{71} \\
\midrule
\midrule
Tokens de texto (BERT)
  & \dims{(B, 100)}
  & \texttt{bert-base-uncased} + padding
  & \fline{datasets/bert_tokenizer.py}{17} \\
\midrule
Embedding
  & \dims{(B, 100, 512)}
  & \texttt{nn.Embedding(29610, 512)}
  & \fline{model/language_network.py}{22} \\
\midrule
BiGRU
  & \dims{(B, 100, 2048)}
  & bidir, $2 \times 1024$ ocultas
  & \fline{model/language_network.py}{24} \\
\midrule
Suma fwd + bwd / 2
  & \dims{(B, 100, 1024)}
  & reduce a 1024 canales
  & \fline{model/language_network.py}{68} \\
\midrule
Permute + unsqueeze
  & \dims{(B, 1024, 100, 1)}
  & prepara para DSC
  & \fline{model/language_network.py}{69} \\
\midrule
Max sobre dim 2
  & \dims{(B, 1024, 1, 1)}
  & pool temporal
  & \fline{model/textreidnet.py}{87} \\
\midrule
DSC 1024 $\to$ 1024
  & \dims{(B, 1024, 1, 1)}
  & conv. separable en profundidad
  & \fline{model/textreidnet.py}{88} \\
\midrule
Textual embedding final
  & \dims{(B, 1024, 1)}
  & \texttt{squeeze(-1)}
  & \fline{model/textreidnet.py}{88} \\
\bottomrule
\end{tabular}
\end{table}

\section{Por qué estas elecciones}

\subsection*{¿Por qué EfficientNet-B0?}

EfficientNet-B0~\cite{tan2019efficientnet} introduce el concepto de
\emph{compound scaling}: en vez de escalar arbitrariamente
profundidad, ancho o resolución, los escala simultáneamente con un
coeficiente fijo derivado de una búsqueda automática. El resultado
es un modelo con excelente relación accuracy/computo. La variante
B0 es la más pequeña de la familia (\textasciitilde 5,3M parámetros)
y es suficiente para la tarea de extracción de features visuales.

\subsection*{¿Por qué BiGRU en lugar de LSTM o Transformer?}

El paper justifica el uso de \gru{} por su eficiencia computacional:
tiene menos parámetros que LSTM (3 compuertas vs 4) y es más rápido
que un Transformer para secuencias cortas (100 tokens). La
bidireccionalidad le permite capturar el contexto anterior y
posterior, lo que mejora la comprensión de la descripción.

\subsection*{¿Por qué convolución separable en profundidad (DSC)?}

Una convolución separable en profundidad~\cite{howard2017mobilenets}
factoriza una convolución estándar en dos operaciones:

\begin{enumerate}
  \item \textbf{Depthwise}: un filtro $k \times k$ por canal de
        entrada.
  \item \textbf{Pointwise}: una convolución $1 \times 1$ para
        mezclar canales.
\end{enumerate}

\noindent Reduce drásticamente los parámetros y las
multiplicaciones, con una pérdida de accuracy mínima. En este
proyecto se usa \emph{dos} veces: para reducir 1280 $\to$ 1024
canales visuales y para proyectar al espacio conjunto.

\subsection*{¿Por qué 1024 como dimensión de embedding?}

1024 es lo suficientemente grande para capturar la semántica de
las dos modalidades, pero suficientemente pequeño para que la
similitud del coseno sea estable y el entrenamiento sea rápido.
Coincide con el número de unidades ocultas de la BiGRU.

\section{Lo que cambia en este proyecto}

El proyecto \texttt{textreid-train} implementa únicamente el
camino de identificación (\emph{TextReIDNet}) sin el detector
R-MHParsNet. Esto significa:

\begin{itemize}
  \item No se detecta automáticamente a las personas en imágenes
        con múltiples personas: la entrada debe ser una imagen con
        una sola persona (post-recortada o de datasets tipo
        CUHK-PEDES).
  \item Se conserva la \emph{colocación sobre canvas negro}
        (\fline{datasets/bases.py}{43}) porque el paper entrena
        así (las imágenes de CUHK-PEDES son menores que $512
        \times 512$ y se colocan aleatoriamente sobre un lienzo
        negro para preservar aspect ratio).
  \item El tamaño de canvas es $512 \times 512$ en el código
        (no $384 \times 128$ como en el paper). Esto es un punto
        a tener en cuenta si se quiere reproducir exactamente los
        números del paper.
\end{itemize}

Estos detalles se discuten con más profundidad en los capítulos
siguientes.

% ============================================================
% Capítulo 3 — Sistema descentralizado
% ============================================================
\chapter{Sistema descentralizado multicámara}
\label{cap:sistema-descentralizado}

Este capítulo describe el sistema completo de vigilancia en el que
se enmarca U-TextReIDNet. \textbf{No hay código en este proyecto
que lo implemente}, pero es útil entenderlo para saber qué parte
del paper es la que realmente vive en \texttt{textreid-train}.

\section{Arquitectura del sistema}

El sistema propuesto en \file{docs/paper.md} (§\textsc{Iii}-B)
tiene tres capas:

\begin{enumerate}
  \item \textbf{Nodos de cámara} ($d_1, \ldots, d_n$): capturan
        vídeo, mantienen galerías locales y ejecutan
        U-TextReIDNet para hacer inferencia local.
  \item \textbf{Infraestructura de comunicaciones}: red entre
        nodos y estación de control.
  \item \textbf{Estación de control}: servidor de aplicaciones
        (con U-TextReIDNet también) + aplicación web
        (\emph{Flask}).
\end{enumerate}

\begin{figure}[H]
\centering
\begin{tikzpicture}[
    font=\sffamily\small,
    node distance=0.5cm and 0.4cm,
    every node/.style={align=center},
    cam/.style={
        rectangle, draw=black, thick, rounded corners=2pt,
        minimum height=1.2cm, minimum width=1.8cm,
        fill=yellow!15
    },
    serv/.style={
        rectangle, draw=black, thick, rounded corners=2pt,
        minimum height=1.5cm, minimum width=3cm,
        fill=red!15
    },
    cli/.style={
        rectangle, draw=black, thick, rounded corners=2pt,
        minimum height=1.2cm, minimum width=2cm,
        fill=cyan!15
    },
    arr/.style={
        -{Stealth[length=2.5mm,width=2mm]}, thick
    }
]

% Nodos de cámara
\node[cam] (d1) {$d_1$\\\footnotesize Jetson};
\node[cam, right=0.5cm of d1] (d2) {$d_2$\\\footnotesize Jetson};
\node[cam, right=0.5cm of d2] (dn) {$d_n$\\\footnotesize Jetson};
\node[right=0.1cm of dn] {\ldots};

% Estación de control
\node[serv, below=2cm of d2] (srv) {\textbf{Servidor de apps}\\CPU i7 + GTX1050TI};

% Cliente
\node[cli, below=1.5cm of srv] (cli) {\textbf{Cliente}\\Navegador web};

% Conexiones
\draw[arr] (d1) -- (srv);
\draw[arr] (d2) -- (srv);
\draw[arr] (dn) -- (srv);
\draw[arr, dashed] (srv) -- (cli) node[midway, right] {\footnotesize Flask};

\end{tikzpicture}
\caption{Topología del sistema descentralizado. Cada cámara corre
U-TextReIDNet localmente; sólo transmite vídeo cuando su puntuación
supera un umbral y es la mejor de todas.}
\label{fig:topologia}
\end{figure}

\section{Flujo operativo}

El paper describe cinco fases (\file{docs/paper.md}, líneas 185--195):

\subsection*{a) Inicialización y solicitud}

Al encenderse, cada cámara se conecta al servidor de aplicaciones
y reporta su estado (\emph{online-idle}, \emph{online-búsqueda},
\emph{transmitiendo}, \emph{offline}). El usuario inicia o detiene
búsquedas desde la aplicación.

\subsection*{b) Captura de fotogramas y galería local}

Cuando una cámara recibe la orden de búsqueda, captura vídeo y
usa R-MHParsNet para generar una galería local de instancias
humanas detectadas en el fotograma actual. \emph{Cada fotograma
se trata como datos nuevos}: no hay relación temporal entre
galerías consecutivas.

\subsection*{c) Inferencia local}

Cada cámara calcula la similitud coseno entre cada imagen de su
galería local y la descripción textual proporcionada por el
usuario. Selecciona la imagen mejor clasificada y, si su
puntuación supera un umbral $\theta$, la procesa para incluir
un bounding box alrededor de la persona objetivo.

\subsection*{d) Identificación global}

Cada segundo, todas las cámaras cuyas puntuaciones superen
$\theta$ envían sus puntuaciones al servidor. El servidor
selecciona la puntuación más alta y notifica a esa cámara para
que transmita vídeo.

\subsection*{e) Finalización}

El usuario debe detener una búsqueda antes de iniciar otra.

\section{Implementación de la estación de control}

La aplicación de usuario es una webapp \emph{Flask}
\cite{grinberg2018flask} que corre en el servidor de aplicaciones.
Se divide en dos subsistemas:

\begin{itemize}
  \item \textbf{Subsistema de comunicación}: usa
        \href{https://zeromq.org/}{ZMQ} para mensajería
        asíncrona entre el servidor y los nodos, y
        \href{https://mqtt.org/}{MQTT} para publicar
        resultados de búsqueda desde los nodos al servidor.
  \item \textbf{Subsistema de vídeo}: usa
        \href{https://github.com/jeffbass/imagezmq}{ImageZMQ}
        para gestionar las transmisiones de vídeo entrantes
        cuando se detecta a la persona objetivo.
\end{itemize}

Estos subsistemas corren en hilos separados para que la UI no se
bloquee durante el procesamiento.

\section{Por qué este proyecto no incluye el sistema}

El repositorio \texttt{textreid-train} es una versión
\emph{training-only} del proyecto original. El README lo explica:

\begin{quote}
El proyecto original es un sistema completo de vigilancia
multicámara descentralizada con scripts de inferencia, una webapp
Flask, MySQL/MQTT y detección humana. \emph{Esta versión quita
todo eso} -- sólo tiene lo necesario para \textbf{entrenar} y
\textbf{evaluar}.
\end{quote}

\noindent Las piezas eliminadas son:

\begin{itemize}
  \item \file{model/unity.py}: wrapper de detección + re-ID
        (innecesario para entrenamiento puro).
  \item \file{model/human_detection_network.py}: detector
        (R-MHParsNet no se entrena aquí).
  \item \file{model/feature_pyramid_network.py}: usado solo por
        el detector.
  \item \file{model/mask_head.py}, \file{parsing_head.py}:
        cabezas del detector.
  \item \file{surveillance_application/}: la webapp Flask.
  \item \file{demo/}: demo web.
  \item \file{inference/}: scripts de inferencia sobre
        imágenes/vídeos.
\end{itemize}

Para desplegar el sistema completo habría que integrar este modelo
con un detector (YOLOv8, por ejemplo) y la infraestructura
descrita arriba. Eso queda fuera del alcance del libro.

\section{Rendimiento reportado en el paper}

El paper mide tiempos de inferencia en dos plataformas:

\begin{table}[H]
\centering
\caption{Tiempos de inferencia reportados en el paper (media sobre
20 iteraciones).}
\label{tab:tiempos-paper}
\begin{tabular}{lrr}
\toprule
\textbf{Modelo} & \textbf{PC (CPU+GPU)} & \textbf{Jetson Nano} \\
\midrule
MFPE        & 0{,}69204 s & --- \\
TextReIDNet & 0{,}00477 s & --- \\
U-TextReIDNet & 0{,}04995 s & 0{,}75013 s \\
\bottomrule
\end{tabular}
\end{table}

\noindent Para una galería de 10 personas:

\begin{table}[H]
\centering
\caption{Tiempos de inferencia para 10 imágenes de galería.}
\begin{tabular}{lrr}
\toprule
\textbf{Modelo} & \textbf{PC} & \textbf{Jetson Nano} \\
\midrule
MFPE          & 2{,}46102 s & --- \\
TextReIDNet   & 1{,}28512 s & --- \\
U-TextReIDNet & 0{,}09525 s & 0{,}75013 s \\
\bottomrule
\end{tabular}
\end{table}

\noindent La conclusión clave es que U-TextReIDNet puede procesar
galerías pequeñas en tiempo real incluso en hardware embebido, y
su tiempo de inferencia escala linealmente con el número de
personas ($O(N \cdot n)$ donde $N$ es la densidad y $n$ el
número de características).

\section{Lo que este proyecto sí hace}

Aunque no incluye el sistema distribuido, el modelo entrenado en
\texttt{textreid-train} \emph{es} el mismo TextReIDNet que se usa
en el paper como submódulo de U-TextReIDNet. Por tanto, un
checkpoint producido aquí puede integrarse en el sistema
completo con sólo cambiar la arquitectura de detección.

% ============================================================
% Capítulo 4 — Dataset CUHK-PEDES
% ============================================================
\chapter{El dataset CUHK-PEDES}
\label{cap:dataset}

Este capítulo describe el dataset con el que se entrena y evalúa
el modelo: qué contiene, qué formato tiene en disco, cómo lo
carga el proyecto y qué decisiones de preprocesamiento se
aplican.

\section{El dataset en cifras}

\textbf{CUHK-PEDES} (\emph{Person Description in the Wild})
\cite{li2017personsearch} es el benchmark estándar para
\reid{} por texto. La versión \emph{original} tiene 40\,206
imágenes y 80\,440 descripciones (2 por imagen). Este proyecto
soporta dos variantes:

\begin{table}[H]
\centering
\caption{Diferencias entre la versión original y la versión
HuggingFace de CUHK-PEDES.}
\label{tab:cuhk-versions}
\begin{tabularx}{\textwidth}{lXX}
\toprule
 & \textbf{Original (paper)} & \textbf{HuggingFace} \\
\midrule
Total imágenes
  & 40\,206
  & 34\,042 \footnotesize(este proyecto) \\
Captions por imagen
  & 2
  & \textasciitilde 7 (1 por fila) \\
Identidades (pids)
  & 13\,003
  & derivadas del hash de imagen \\
Splits
  & train / val / test
  & derivados: 75\% / 12,5\% / 12,5\% \\
Resolución
  & variable (\textasciitilde 256×128+)
  & 128×128 thumbnails \\
Acceso
  & Email a Tong Xiao
  & Inmediato desde HuggingFace \\
Top-1 esperado
  & \textasciitilde 54\% (paper)
  & \textasciitilde 45-52\% \\
\bottomrule
\end{tabularx}
\end{table}

\noindent\textbf{Recomendación}: si se pueden obtener los datos
originales por email (gratis para investigación), se obtienen
mejores números. Si no, la versión HF funciona perfectamente
para validar el pipeline.

\section{Formato en disco}

Después de ejecutar \file{scripts/download_hf_dataset.py}, la
versión HF queda en \file{data/CUHK-PEDES-HF/} con esta
estructura:

\begin{verbatim}
data/CUHK-PEDES-HF/
|---- raw.parquet               # snapshot de HuggingFace (550 MB)
|---- reid_raw.json             # anotaciones en formato "estándar"
|__-- imgs/
    |-- imgs/all/
        |-- 7138159738431296927.jpg
        |-- 3825926175112857003.jpg
        `-- ... (34042 imágenes en total)
\end{verbatim}

\subsection{El archivo \texttt{reid\_raw.json}}

Es la \emph{fuente de verdad} del dataset. Cada elemento es un
diccionario con estos campos:

\begin{minted}[basicstyle=\ttfamily\small, frame=single]{python}
{
  "file_path": "all/7138159738431296927.jpg",
  "captions": [
    "A pedestrian with dark hair is wearing red and white shoes, "
    "a black hooded sweatshirt, and black pants.",
    "The person has short black hair and is wearing black pants, "
    "a long sleeve black top, and red sneakers.",
    ...
  ],
  "split": "train",
  "processed_tokens": [],
  "id": 0
}
\end{minted}

\noindent Los campos obligatorios son sólo \texttt{file\_path} y
\texttt{captions}. Los demás se completan en tiempo de carga.

\section{Cómo se cargan los datos}

El proyecto tiene dos clases de carga según la fuente:

\begin{itemize}
  \item \file{datasets/cuhkpedes.py}: \texttt{CUHKPEDES} para
        datos originales.
  \item \file{datasets/cuhkpedes_hf.py}: \texttt{CUHKPEDESHF}
        para datos de HuggingFace.
\end{itemize}

Ambas clases exponen la misma interfaz:

\begin{minted}[basicstyle=\ttfamily\small, frame=single]{python}
ds.train        # lista de (pid, image_id, img_path, caption)
ds.train_id_container    # set de pids en train
ds.val          # dict con image_pids, img_paths, caption_pids, captions
ds.test         # dict idem para test
\end{minted}

\subsection{El loader HF: \texttt{CUHKPEDESHF}}

Esta clase (\file{datasets/cuhkpedes_hf.py}) tiene tres pasos
clave:

\begin{enumerate}
  \item \textbf{Asignación de pids}: como la versión HF no
        incluye IDs de persona, los deriva del \emph{hash de la
        imagen}. Las imágenes idénticas (mismas captions)
        comparten el mismo pid.
  \item \textbf{División train/val/test}: por proporción de IDs:
        75\% / 12,5\% / 12,5\%. No usa el campo
        \texttt{split} del JSON.
  \item \textbf{Procesamiento de anotaciones}: igual que la
        versión original.
\end{enumerate}

\noindent El código clave está en \fline{datasets/cuhkpedes_hf.py}{62-71}:

\begin{minted}[numbers, firstnumber=62, linerange=62-71]{python}
def _build_pid_mapping(self):
    """Group entries by image_path (which is the identity key)."""
    self.path_to_id = {}
    next_id = 0
    for anno in self.annos:
        fp = anno['file_path']
        if fp not in self.path_to_id:
            self.path_to_id[fp] = next_id
            next_id += 1
    self.num_persons = next_id
\end{minted}

\subsection{El loader original: \texttt{CUHKPEDES}}

La clase \file{datasets/cuhkpedes.py} asume que las anotaciones
vienen con \texttt{split} explícito (train/val/test) y con
\texttt{id} entre 1 y 13\,003. Internamente resta 1 al id para
que empiece en 0 (más conveniente para PyTorch). El código
relevante está en \fline{datasets/cuhkpedes.py}{75-115}.

\section{El dataloader unificado}

\file{datasets/cuhkpedes_dataloader.py} es el punto de entrada
que decide qué loader usar y construye los \texttt{DataLoader}s
de PyTorch. Su función principal es
\texttt{build\_cuhkpedes\_dataloader(config)}.

\begin{minted}[numbers, firstnumber=22, linerange=22-86]{python}
def build_cuhkpedes_dataloader(config: dict = None):
    if config.dataset_source == 'huggingface':
        dataset_object = CUHKPEDESHF(config)
    else:
        dataset_object = CUHKPEDES(config)

    train_transform = T.Compose([
        T.ToTensor(),
        T.Normalize(mean=config.MHPV2_means, std=config.MHPV2_stds)
    ])
    inference_transform = T.Compose([
        T.ToTensor(),
        T.Normalize(mean=config.MHPV2_means, std=config.MHPV2_stds)
    ])
    train_num_classes = len(dataset_object.train_id_container)

    train_set = ImageTextDataset(
        dataset_object.train,
        train_transform,
        tokenizer_type=config.tokenizer_type,
        tokens_length_max=config.tokens_length_max
    )

    train_data_loader = DataLoader(
        train_set,
        batch_size=config.batch_size,
        shuffle=True,
        num_workers=config.num_workers,
        collate_fn=collate
    )
    ...
    return train_data_loader, inference_img_loader,
           inference_txt_loader, train_num_classes
\end{minted}

\noindent\textbf{Nota importante}: las transformaciones de
imagen \emph{sólo} convierten a tensor y normalizan. El resize a
$384 \times 128$ que el paper describe \textbf{no está}. Las
imágenes se quedan en su tamaño nativo ($512 \times 512$ después
del \texttt{place\_image\_on\_canvas}).

\section{La clase \texttt{ImageTextDataset}}

Es el \texttt{Dataset} de PyTorch que produce un batch mixto
imagen+texto. Vive en \file{datasets/bases.py}. Para cada índice:

\begin{minted}[numbers, firstnumber=124, linerange=124-151]{python}
def __getitem__(self, index):
    pid, image_id, img_path, caption = self.dataset[index]
    pid = torch.from_numpy(np.array([pid])).long()
    img = read_image(img_path)
    img = place_image_on_canvas(img)
    original_img = img.copy()
    if self.transform is not None:
        img = self.transform(img)
    # ... guarda también la imagen original sin normalizar
    tokens = self.tokenizer(caption)
    token_ids, orig_token_length = pad_tokens(tokens,
                                              self.tokens_length_max)

    ret = {'pids': pid,
           'image_ids': image_id,
           'img_paths': img_path,
           'preprocessed_images': img,
           'original_images': original_img,
           'token_ids': token_ids.to(torch.long),
           'orig_token_lengths': orig_token_length,
           'captions': caption}
    return ret
\end{minted}

\noindent Cada item del dataset devuelve un diccionario con 8
claves. El \texttt{collate} (\fline{utils/miscellaneous_utils.py}{189})
se encarga de apilar tensores y dejar las cadenas como listas.

\section{Preprocesamiento de imagen: \texttt{place\_image\_on\_canvas}}

Función definida en \file{datasets/bases.py}:

\begin{minted}[numbers, firstnumber=43, linerange=43-83]{python}
def place_image_on_canvas(image, new_size=(512, 512),
                          placement_strategy='random'):
    """Plaza una imagen en un canvas negro preservando aspect ratio."""
    if image.height > 512:
        aspect_ratio = image.width / image.height
        new_height, _ = new_size
        new_width = int(aspect_ratio * new_height)
        image = image.resize((new_width, new_height), Image.ANTIALIAS)

    canvas = Image.new('RGB', new_size, (0, 0, 0))

    if placement_strategy == 'center':
        x = (new_size[0] - image.width) // 2
        y = (new_size[1] - image.height) // 2
    elif placement_strategy == 'random':
        x = np.random.randint(0, max(1, new_size[0] - image.width))
        y = np.random.randint(0, max(1, new_size[1] - image.height))
    else:
        raise ValueError(...)

    canvas.paste(image, (x, y))
    return canvas
\end{minted}

\noindent\textbf{¿Por qué?} Las imágenes de CUHK-PEDES son
chicas ($\textasciitilde 128 \times 128$ en HF,
$\textasciitilde 256 \times 128$ en original). EfficientNet-B0
espera imágenes relativamente grandes ($224 \times 224$ o
más). Para no distorsionar el aspect ratio, se pega la imagen en
una posición aleatoria sobre un canvas negro de $512 \times 512$.
Esto introduce variación posicional como augmentation implícita.

\section{Preprocesamiento de texto}

\subsection*{Tokenización}

Hay tres tokenizers disponibles; el que se usa por defecto es
\textbf{BERT} (\texttt{bert-base-uncased}).

\begin{minted}[numbers, firstnumber=10, linerange=10-19]{python}
class BERTTokenizer(object):
    def __init__(self):
        super(BERTTokenizer, self).__init__()
        _ , tokenizer_class, pretrained_weights = (
            ppb.BertModel, ppb.BertTokenizer, 'bert-base-uncased')
        self.tokenizer = tokenizer_class.from_pretrained(
            pretrained_weights)

    def __call__(self, text: str = None) -> torch.LongTensor:
        tokens = self.tokenizer.encode(text)
        return tokens
\end{minted}

\noindent\textbf{Otros tokenizers} disponibles en
\file{datasets/simple_tokenizer.py} y
\file{datasets/tiktoken_tokenizer.py}, pero el config
(\fline{config.py}{41}) fija \texttt{tokenizer\_type='bert'} por
defecto.

\subsection*{Padding}

La función \texttt{pad\_tokens}
(\fline{utils/miscellaneous_utils.py}{86-107}) recorta o
padea a longitud fija:

\begin{minted}[numbers, firstnumber=86, linerange=86-107]{python}
def pad_tokens(tokens, tokens_length_max):
    tokens_length = len(tokens)
    tokens = torch.from_numpy(np.array(tokens)).view(-1)
                                  .contiguous().float()
    if tokens_length < tokens_length_max:
        zero_padding = torch.zeros(tokens_length_max - tokens_length)
        tokens = torch.cat([tokens, zero_padding], 0)
    else:
        tokens = tokens[:tokens_length_max]
        tokens_length = tokens_length_max
    return tokens, tokens_length
\end{minted}

\noindent\textbf{Tokens de padding = 0}. Esto es importante
porque el \texttt{nn.Embedding} del modelo tiene
\texttt{padding\_idx=0} (\fline{model/language_network.py}{22}),
lo que garantiza que el padding no afecte al embedding.

\section{Normalización de imagen}

El dataloader aplica \texttt{T.Normalize} con medias y desviaciones
de MHPV2 (no de ImageNet):

\begin{minted}[basicstyle=\ttfamily\small, frame=single]{python}
mean = [0.3555248975753784, 0.3295726776123047, 0.3115333914756775]
std  = [0.3482806384563446, 0.3339986503124237, 0.3296058773994446]
\end{minted}

\noindent definidos en \fline{config.py}{50-51}. \textbf{Esto es
un poco raro}: el paper usaría medias de ImageNet ($[0.485, 0.456,
0.406]$) que son las guardadas en \texttt{CUHKPEDES\_means}
(\fline{config.py}{52}). El código actual usa MHPV2 (que es
\emph{otro} dataset). Es un detalle a tener en cuenta al
reproducir resultados.

\section{El cuaderno de validación}

En \file{notebooks/01\_lab\_validate\_dataset.ipynb} hay un
notebook de Jupyter que verifica:

\begin{enumerate}
  \item Que el JSON existe y tiene los campos correctos.
  \item Que las imágenes referenciadas existen y son válidas.
  \item La distribución de longitud de captions.
  \item La distribución de splits.
  \item Que un forward pass del modelo funciona.
  \item Las pérdidas dan valores finitos.
  \item Una estimación de VRAM para elegir el \texttt{batch\_size}.
\end{enumerate}

\noindent Conviene correr este notebook \emph{antes} del primer
entrenamiento, especialmente si se usa una GPU diferente a la
del paper.

\section{Dataset augmentation}

Este proyecto \textbf{no aplica augmentation} explícita de
imagen más allá de:

\begin{itemize}
  \item Posición aleatoria en el canvas (\texttt{placement\_strategy='random'}).
  \item Normalización con medias MHPV2.
\end{itemize}

\noindent El paper original describe un pipeline más rico
(\fline{docs/paper.md}{245}):
\textit{redimensionamiento, recorte aleatorio y normalización}.
Este proyecto lo omite; podría añadirse en
\texttt{get\_transform} (\fline{utils/miscellaneous\_utils.py}{210}).

\section{Resumen del pipeline de datos}

\begin{figure}[H]
\centering
\begin{tikzpicture}[
    font=\sffamily\small,
    node distance=0.6cm,
    box/.style={
        rectangle, draw=black, thick, rounded corners=2pt,
        minimum height=0.9cm, minimum width=2.2cm, align=center,
        fill=blue!10
    },
    arr/.style={-{Stealth[length=2.5mm,width=2mm]}, thick}
]

\node[box] (json)   {reid\_raw.json};
\node[box, right=of json] (loader) {CUHKPEDES\\(HF)};
\node[box, right=of loader] (ds) {ImageTextDataset};
\node[box, right=of ds] (dl)  {DataLoader\\(collate)};
\node[box, right=of dl]  (out) {Batch\\\{preprocessed\_images,\\token\_ids,\\pids, ...\}};

\draw[arr] (json)   -- (loader);
\draw[arr] (loader) -- (ds);
\draw[arr] (ds)     -- (dl);
\draw[arr] (dl)     -- (out);

\node[below=0.2cm of loader, font=\scriptsize] {image\_pids, captions};
\node[below=0.2cm of ds,     font=\scriptsize] {place on canvas};
\node[below=0.2cm of dl,     font=\scriptsize] {shuffle, workers};
\end{tikzpicture}
\caption{Pipeline de datos desde el JSON hasta el batch que
consume \texttt{train.py}.}
\label{fig:pipeline-datos}
\end{figure}

% ============================================================
% Capítulo 5 — Configuración
% ============================================================
\chapter{Configuración del proyecto}
\label{cap:configuracion}

Todo el comportamiento del proyecto está controlado por
\file{config.py}. Este archivo define un diccionario de
configuración (\texttt{DotDict}) que se pasa a cada módulo y
permite modificar el experimento sin tocar el código.

\section{La clase \texttt{DotDict}}

\begin{minted}[numbers, firstnumber=14, linerange=14-17]{python}
class DotDict(dict):
    __getattr__ = dict.get
    __setattr__ = dict.__setitem__
    __delattr__ = dict.__delitem__
\end{minted}

\noindent Es un \texttt{dict} con acceso por atributo:
\texttt{configs.lr} es equivalente a \texttt{configs['lr']}.
Esto permite usar tanto \texttt{config.feature\_length} como
\texttt{config['feature\_length']} en el resto del código.

\section{Función principal: \texttt{sys\_configuration}}

Es la única función exportada. Acepta tres parámetros y devuelve
un \texttt{DotDict}:

\begin{minted}[numbers, firstnumber=20, linerange=20-31]{python}
def sys_configuration(platform_name: str = platform.node(),
                     dataset_name: str = "CUHK-PEDES",
                     dataset_source: str = "huggingface",
                     show_configs: bool = False) -> DotDict:
    """
    Args.:
        platform_name: hostname of the machine
        dataset_name:  "CUHK-PEDES" (only one supported for now)
        dataset_source: "huggingface" or "original"
    """
\end{minted}

\noindent Los argumentos se pueden sobreescribir desde CLI en
\texttt{train.py} (\fline{train.py}{56}).

\section{Tabla maestra de parámetros}

\subsection*{Modelo de lenguaje}

\begin{longtable}{p{3.5cm}p{2.5cm}p{8cm}}
\toprule
\textbf{Parámetro} & \textbf{Valor} & \textbf{Efecto / ubicación} \\
\midrule
\endhead
\texttt{dropout\_rate}
  & 0.30
  & Dropout tras el embedding textual. \\
  & & \fline{config.py}{37}, \fline{model/language_network.py}{23}. \\
\midrule
\texttt{num\_layers}
  & 1
  & Número de capas del \gru{} bidireccional. \\
  & & \fline{config.py}{38}, \fline{model/language_network.py}{24}. \\
\midrule
\texttt{tokens\_length\_max}
  & 100
  & Longitud fija (con padding) de cada caption tokenizada. \\
  & & \fline{config.py}{39}, \fline{datasets/bases.py}{96}. \\
\midrule
\texttt{feature\_length}
  & 1024
  & Dimensión del espacio de embedding conjunto. \\
  & & \fline{config.py}{40}, \fline{model/textreidnet.py}{43}. \\
\midrule
\texttt{tokenizer\_type}
  & 'bert'
  & Tipo de tokenizer. Otros: 'simple\_tokenizer', 'tiktoken\_*'. \\
  & & \fline{config.py}{41}, \fline{datasets/bases.py}{97}. \\
\midrule
\texttt{vocab\_size}
  & 29\,610
  & 29\,609 (BERT) + 1 (padding). Padding idx = 0. \\
  & & \fline{config.py}{42}, \fline{model/language_network.py}{22}. \\
\midrule
\texttt{embedding\_dim}
  & 512
  & Dimensión del embedding textual. \\
  & & \fline{config.py}{43}, \fline{model/language_network.py}{22}. \\
\bottomrule
\caption{Parámetros del modelo de lenguaje.}
\label{tab:cfg-lang}
\end{longtable}

\subsection*{Modelo visual}

\begin{longtable}{p{3.5cm}p{2.5cm}p{8cm}}
\toprule
\textbf{Parámetro} & \textbf{Valor} & \textbf{Efecto / ubicación} \\
\midrule
\endhead
\texttt{CUHKPEDES\_image\_size}
  & (384, 128)
  & Tamaño \emph{declarado} en el paper. \textbf{No se usa}
        en el dataloader actual. \\
  & & \fline{config.py}{48}. \\
\midrule
\texttt{MHPV2\_image\_size}
  & (512, 512)
  & Tamaño del canvas en el que se pegan las imágenes. \\
  & & \fline{config.py}{49}, \fline{datasets/bases.py}{58}. \\
\midrule
\texttt{MHPV2\_means}
  & [0.355, 0.330, 0.311]
  & Medias para \texttt{T.Normalize} en el dataloader. \\
  & & \fline{config.py}{50}. \\
\midrule
\texttt{MHPV2\_stds}
  & [0.348, 0.334, 0.330]
  & Desviaciones para \texttt{T.Normalize}. \\
  & & \fline{config.py}{51}. \\
\midrule
\texttt{CUHKPEDES\_means}
  & [0.485, 0.456, 0.406]
  & Medias de ImageNet (no usadas en este proyecto). \\
  & & \fline{config.py}{52}. \\
\midrule
\texttt{CUHKPEDES\_stds}
  & [0.229, 0.224, 0.225]
  & Desviaciones de ImageNet (no usadas). \\
  & & \fline{config.py}{53}. \\
\bottomrule
\caption{Parámetros del modelo visual.}
\label{tab:cfg-vis}
\end{longtable}

\subsection*{Entrenamiento}

\begin{longtable}{p{3.5cm}p{2.5cm}p{8cm}}
\toprule
\textbf{Parámetro} & \textbf{Valor} & \textbf{Efecto / ubicación} \\
\midrule
\endhead
\texttt{epoch}
  & 60
  & Número de epochs total. \\
  & & \fline{config.py}{71}, \fline{train.py}{113}. \\
\midrule
\texttt{adam\_alpha}
  & 0.90
  & $\beta_1$ de AdamW. \\
  & & \fline{config.py}{72}, \fline{train.py}{98}. \\
\midrule
\texttt{adam\_beta}
  & 0.999
  & $\beta_2$ de AdamW. \\
  & & \fline{config.py}{73}, \fline{train.py}{98}. \\
\midrule
\texttt{epoch\_decay}
  & [20, 40]
  & Épocas donde el LR scheduler decae el learning rate. \\
  & & \fline{config.py}{74}, \fline{train.py}{99}. \\
\midrule
\texttt{lr}
  & 0.001
  & Learning rate inicial. \\
  & & \fline{config.py}{75}, \fline{train.py}{98}. \\
\midrule
\texttt{patience}
  & 3
  & (No usado actualmente; sería early stopping). \\
  & & \fline{config.py}{76}. \\
\midrule
\texttt{val\_dataset}
  & 'test'
  & Split usado durante validación. \\
  & & \fline{config.py}{77}. \\
\midrule
\texttt{seed}
  & 3407
  & Semilla para reproducibilidad. \\
  & & \fline{config.py}{79}, \fline{train.py}{71}. \\
\midrule
\texttt{batch\_size}
  & 8 (default)
  & Tamaño de batch (4 en Jetson). \\
  & & \fline{config.py}{123,126,129}. \\
\midrule
\texttt{num\_workers}
  & 4 (default)
  & Hilos del DataLoader. \\
  & & \fline{config.py}{122,125,128}. \\
\bottomrule
\caption{Parámetros de entrenamiento.}
\label{tab:cfg-train}
\end{longtable}

\subsection*{Pérdidas}

\begin{longtable}{p{3.5cm}p{2.5cm}p{8cm}}
\toprule
\textbf{Parámetro} & \textbf{Valor} & \textbf{Efecto / ubicación} \\
\midrule
\endhead
\texttt{margin}
  & 0.5
  & Margen del triplet-style RankingLoss. \\
  & & \fline{config.py}{86}, \fline{evaluation/ranking_loss.py}{41}. \\
\midrule
\texttt{ranking\_loss\_alpha}
  & 1.0
  & Peso $\alpha$ de la RankingLoss en la pérdida total. \\
  & & \fline{config.py}{87}, \fline{train.py}{145}. \\
\midrule
\texttt{identity\_loss\_beta}
  & 1.0
  & Peso $\beta$ de la IdentityLoss en la pérdida total. \\
  & & \fline{config.py}{88}, \fline{train.py}{146}. \\
\bottomrule
\caption{Parámetros de funciones de pérdida.}
\label{tab:cfg-loss}
\end{longtable}

\subsection*{Paths y logging}

\begin{longtable}{p{3.5cm}p{5cm}p{5.5cm}}
\toprule
\textbf{Parámetro} & \textbf{Valor} & \textbf{Ubicación} \\
\midrule
\endhead
\texttt{train\_log\_path}
  & \file{logs/train.log}
  & \fline{config.py}{93} \\
\texttt{test\_log\_path}
  & \file{logs/test.log}
  & \fline{config.py}{94} \\
\texttt{model\_save\_path}
  & \file{data/checkpoints}
  & \fline{config.py}{96} \\
\texttt{plot\_save\_path}
  & \file{data/plots}
  & \fline{config.py}{97} \\
\texttt{dataset\_path} (HF)
  & \file{data/CUHK-PEDES-HF}
  & \fline{config.py}{108} \\
\texttt{dataset\_path} (orig)
  & \file{data/CUHK-PEDES}
  & \fline{config.py}{112} \\
\bottomrule
\caption{Paths del proyecto.}
\label{tab:cfg-paths}
\end{longtable}

\section{Override por línea de comandos}

\file{train.py} permite sobreescribir los parámetros más
importantes:

\begin{minted}[basicstyle=\ttfamily\small, frame=single]{python}
python train.py \
    --dataset_source huggingface \
    --epochs 60 \
    --batch_size 8 \
    --lr 0.001 \
    --seed 3407 \
    --output_dir ./data/checkpoints
\end{minted}

\noindent Cada flag corresponde a un campo del config que se
sobreescribe justo después de llamar a \texttt{sys\_configuration}
(\fline{train.py}{56-69}).

\section{Configuración por plataforma}

Hay una rama de detección automática de plataforma basada en el
hostname (\fline{config.py}{121-129}):

\begin{minted}[numbers, firstnumber=121, linerange=121-129]{python}
if platform_name in ('PC-SIM', 'deeplearning', 'ultron', 'yrsn'):
    configs['num_workers'] = 4
    configs['batch_size'] = 8
elif platform_name == 'nano':
    configs['num_workers'] = 2
    configs['batch_size'] = 4
else:
    configs['num_workers'] = 4
    configs['batch_size'] = 8
\end{minted}

\noindent Si tu hostname es ``nano'' (Jetson Nano), el batch se
reduce automáticamente a 4. Para cualquier otro caso, usa 8.

\section{Reproducibilidad}

\texttt{set\_seed}
(\fline{utils/miscellaneous_utils.py}{18-30}) fija:

\begin{itemize}
  \item \texttt{PYTHONHASHSEED}.
  \item \texttt{random.seed}, \texttt{numpy.random.seed}.
  \item \texttt{torch.manual\_seed} y \texttt{torch.cuda.manual\_seed}.
  \item \texttt{cudnn.deterministic=True}, \texttt{cudnn.benchmark=False}.
\end{itemize}

\noindent Esto garantiza resultados idénticos entre ejecuciones
en la misma máquina. \textbf{No} garantiza resultados idénticos
entre máquinas distintas.

\section{Ejemplo de inspección}

Para imprimir la configuración completa desde un script:

\begin{minted}[basicstyle=\ttfamily\small, frame=single]{python}
from config import sys_configuration
cfg = sys_configuration(dataset_source='huggingface')
print(cfg.feature_length)    # 1024
print(cfg.batch_size)        # 8
print(cfg.vocab_size)        # 29610
\end{minted}

\noindent Esta es la mejor forma de verificar que la
configuración es la esperada antes de entrenar.

% ============================================================
% Capítulo 6 — Visual Network: conceptos + código
% ============================================================
\chapter{Red visual: EfficientNet-B0 y convolución separable}
\label{cap:visual-network}

Este capítulo combina \textbf{explicaciones conceptuales} de las
piezas usadas en la rama visual con el \textbf{código línea por
línea} del proyecto.

\section{¿Qué es EfficientNet?}

\subsection{El problema}

Hasta 2019, la mayoría de las CNNs se escalaban de forma
\emph{arbitraria}: o se hacían más profundas (más capas), o más
anchas (más canales), o con mayor resolución de entrada. Cada
familia tenía su propia receta y nadie sabía cuál era la mejor.

\subsection{Compound scaling}

\eff{}~\cite{tan2019efficientnet} introdujo una idea elegante:
\textbf{escalar las tres dimensiones simultáneamente} con un
ratio fijo ($\phi$). El paper demostró que esto da mejor
accuracy por FLOP que cualquier escalado aislado.

\begin{equation}
  \begin{aligned}
    \text{depth}   &= \alpha^{\phi} \\
    \text{width}   &= \beta^{\phi} \\
    \text{resolution} &= \gamma^{\phi} \\
    \text{con} \quad & \alpha \cdot \beta^2 \cdot \gamma^2 \approx 2
  \end{aligned}
  \label{eq:compound-scaling}
\end{equation}

\noindent La familia EfficientNet va desde B0 (la más chica) hasta
B7 (la más grande). En este proyecto se usa \textbf{B0} porque es
la única que cabe en hardware modesto con un budget razonable.

\subsection{¿Qué es MBConv?}

MBConv (\emph{Mobile Inverted Bottleneck Convolution}) es el
bloque básico de EfficientNet. Inventado por MobileNetV2 y
mejorado por EfficientNet. Tiene tres etapas:

\begin{enumerate}
  \item \textbf{Expansión}: conv 1x1 que aumenta canales
        ($\times 6$ típicamente, por eso "MBConv6").
  \item \textbf{Depthwise conv}: conv $k \times k$ por canal
        (máscara espacial).
  \item \textbf{Proyección}: conv 1x1 que reduce canales.
\end{enumerate}

\noindent A diferencia del bottleneck clásico de ResNet, MBConv
\emph{expande} primero y \emph{comprime} después ("inverted").

\subsection{Layout de B0}

\begin{table}[H]
\centering
\caption{Las 9 etapas de EfficientNet-B0.}
\label{tab:effnet-stages-concepto}
\begin{tabular}{clrrr}
\toprule
\textbf{Stage} & \textbf{Operador} & \textbf{Resolución} & \textbf{Canales} & \textbf{Layers} \\
\midrule
1  & Conv3x3       & $224 \times 224$ & 32   & 1 \\
2  & MBConv1, k3x3 & $112 \times 112$ & 16   & 1 \\
3  & MBConv6, k3x3 & $112 \times 112$ & 24   & 2 \\
4  & MBConv6, k5x5 & $56 \times 56$   & 40   & 2 \\
5  & MBConv6, k3x3 & $28 \times 28$   & 80   & 3 \\
6  & MBConv6, k5x5 & $14 \times 14$   & 112  & 3 \\
7  & MBConv6, k5x5 & $14 \times 14$   & 192  & 4 \\
8  & MBConv6, k3x3 & $7 \times 7$     & 320  & 1 \\
9  & Conv1x1+Pool  & $7 \times 7$     & 1280 & 1 \\
\bottomrule
\end{tabular}
\end{table}

\noindent\textbf{Cada etapa divide la resolución a la mitad}
(excepto 2 y 5 que mantienen). La \textbf{etapa 9} es la que se usa
en TextReIDNet: produce 1280 canales a $7 \times 7$ (con entrada
$224 \times 224$) o $16 \times 16$ (con entrada $512 \times 512$).

\section{¿Qué es la convolución separable en profundidad (DSC)?}

\subsection{El problema}

Una convolución 2D estándar $k \times k$ sobre una imagen con $C$
canales de entrada y $C'$ canales de salida tiene:
\begin{itemize}
  \item \textbf{Parámetros}: $k \cdot k \cdot C \cdot C'$.
  \item \textbf{FLOPs}: $k^2 \cdot C \cdot C' \cdot H \cdot W$.
\end{itemize}

\noindent Para $k=3$, $C=C'=1280$: \textbf{14,7M} parámetros y
muchos FLOPs.

\subsection{La factorización}

La \emph{depthwise separable convolution} (DSC)~\cite{howard2017mobilenets}
factoriza esa operación en dos:

\begin{enumerate}
  \item \textbf{Depthwise}: una convolución $k \times k$ por canal,
        con $k^2 \cdot C$ parámetros (en vez de $k^2 \cdot C \cdot C'$).
  \item \textbf{Pointwise}: una convolución $1 \times 1$ para mezclar
        canales, con $C \cdot C'$ parámetros.
\end{enumerate}

\noindent\textbf{Resultado}: $k^2 \cdot C + C \cdot C'$ parámetros,
una reducción de $\sim k^2$ veces. Para $k=3$: $\sim 9\times$.

\subsection{Visualización}

\begin{figure}[H]
\centering
\begin{tikzpicture}[
    node distance=0.4cm and 0.6cm,
    every node/.style={font=\sffamily\footnotesize, align=center},
    box/.style={
        rectangle, draw=black, thick, rounded corners=2pt,
        minimum height=1.2cm, minimum width=1.6cm, fill=blue!10
    },
    arr/.style={-{Stealth[length=2.5mm,width=2mm]}, thick}
]
% Standard
\node[box] (in1) {Input\\$C$};
\node[box, right=of in1] (k1) {Conv $k{\times}k$\\$C \to C'$};
\node[box, right=of k1] (out1) {Output\\$C'$};
\draw[arr] (in1) -- (k1); \draw[arr] (k1) -- (out1);
\node[above=0.3cm of k1] {Estándar};

% DSC
\node[box, below=2cm of in1] (in2) {Input\\$C$};
\node[box, right=of in2] (d) {Depthwise\\$k{\times}k$\\por canal};
\node[box, right=of d] (p) {Pointwise\\$1{\times}1$\\$C \to C'$};
\node[box, right=of p, minimum width=1.6cm] (out2) {Output\\$C'$};
\draw[arr] (in2) -- (d); \draw[arr] (d) -- (p); \draw[arr] (p) -- (out2);
\node[above=0.3cm of d] {DSC};

\node[below=1.5cm of p, font=\small, text width=12cm, align=center]
    {Parámetros: estándar $= k^2 \cdot C \cdot C'$ \\[2pt]
    DSC $= k^2 \cdot C + C \cdot C'$ \\[-2pt]
    (reducción $\approx k^2$ veces)};
\end{tikzpicture}
\caption{Comparación entre convolución estándar y DSC.}
\label{fig:dsc}
\end{figure}

\section{¿Qué es AdaptiveMaxPool2d?}

Un \emph{max pooling adaptativo} recibe como argumento el tamaño
\emph{deseado de salida}, no el stride ni el kernel. Calcula
automáticamente los hiperparámetros para que el tensor de salida
tenga exactamente esa forma.

\begin{itemize}
  \item \texttt{nn.AdaptiveMaxPool2d((1, 1))}: salida siempre
        $(B, C, 1, 1)$.
  \item Para cada canal $c$, devuelve el valor máximo sobre todas
        las posiciones espaciales.
  \item Equivale a un \emph{GAP} (global max pooling) si el kernel
        efectivo cubre toda la imagen.
\end{itemize}

\noindent\textbf{¿Por qué?} Después del EfficientNet tenemos un
tensor $(B, 1024, 16, 16)$. Necesitamos un vector por imagen. El
max pooling global colapsa $(16, 16) \to (1, 1)$ quedándose con la
activación más fuerte de cada canal.

\section{Código: la clase \texttt{VisualNetwork}}

\begin{minted}[language=Python, numbers, firstnumber=7, linerange=7-39]{python}
class VisualNetwork(nn.Module):
    """Backbone visual: EfficientNet-B0 preentrenado."""

    def __init__(self):
        # Constructor de nn.Module (registra parámetros)
        super(VisualNetwork, self).__init__()

        # Cargar EfficientNet-B0 con pesos preentrenados en ImageNet
        # EfficientNet_B0_Weights.DEFAULT = mejores pesos disponibles
        # (~5.3M parámetros, input esperado 224x224)
        self.model = models.efficientnet_b0(
            weights=EfficientNet_B0_Weights.DEFAULT
        )

        # `model.features` es un nn.Sequential con 9 bloques
        # (ver Tabla 6.1). Vamos a iterar bloque por bloque.
        self.feature_blocks = self.model.features

    def forward(self, x):
        """x: (B, 3, H, W) imagen preprocesada.
           returns: (B, 1280, H/32, W/32) salida etapa 9."""
        # Lista para acumular features de cada bloque
        features = []

        # Iterar sobre los 9 bloques secuenciales de EfficientNet-B0
        for block in self.feature_blocks:
            # Aplicar bloque (conv + BN + activation)
            x = block(x)

            # Guardar feature de este bloque
            features.append(x)

        # Devolver SOLO la salida del último bloque (etapa 9)
        # = 1280 canales, resolución espacial /32
        return features[-1]
\end{minted}

\section{Código: \texttt{DepthwiseSeparableConv}}

\begin{minted}[language=Python, numbers, firstnumber=259, linerange=259-287]{python}
class DepthwiseSeparableConv(nn.Module):
    """Convolución separable en profundidad: depthwise + pointwise."""

    def __init__(self, in_channels=None, out_channels=None,
                 kernel_size=3, bias=True, is_final_conv=False):
        # Constructor de nn.Module
        super(DepthwiseSeparableConv, self).__init__()

        # Depthwise convolution: una convolución k×k por canal
        # groups=in_channels hace que cada filtro actúe sobre UN canal
        # padding=1 mantiene la dimensión espacial
        self.depthwise_convolution = nn.Conv2d(
            in_channels=in_channels,
            out_channels=in_channels,    # mismo número de canales
            kernel_size=kernel_size,
            groups=in_channels,          # ← clave: depthwise
            bias=bias,
            padding=1                    # mantiene H, W
        )

        # Pointwise convolution: 1×1 que mezcla canales
        # in_channels -> out_channels
        self.pointwise_convolution = nn.Conv2d(
            in_channels=in_channels,
            out_channels=out_channels,   # ← dimensión de salida
            kernel_size=1,
            groups=1,                    # convolución normal
            bias=bias
        )

        # Encadenar las dos operaciones
        self.depthwise_separable_conv = nn.Sequential(
            self.depthwise_convolution,
            self.pointwise_convolution
        )

        # RetinaNet-style bias init (no se usa en TextReIDNet)
        self.is_final_conv = is_final_conv
        if bias and is_final_conv:
            self.is_final_conv = True

        # Inicialización Kaiming (para Conv2d y Linear)
        self.depthwise_separable_conv.apply(weights_init_kaiming)

    def forward(self, x):
        return self.depthwise_separable_conv(x)
\end{minted}

\section{Etapas con shapes exactas (input 512×512)}

\begin{table}[H]
\centering
\caption{Etapas con shapes para input $512 \times 512$ en este proyecto.}
\label{tab:effnet-shapes}
\begin{tabular}{clrrr}
\toprule
\textbf{Stage} & \textbf{Operador} & \textbf{Output HxW} & \textbf{Output Ch} & \textbf{Layers} \\
\midrule
0  & Conv3x3       & $256 \times 256$  & 32   & 1 \\
1  & MBConv1, k3x3 & $256 \times 256$  & 16   & 1 \\
2  & MBConv6, k3x3 & $128 \times 128$  & 24   & 2 \\
3  & MBConv6, k5x5 & $64 \times 64$    & 40   & 2 \\
4  & MBConv6, k3x3 & $32 \times 32$    & 80   & 3 \\
5  & MBConv6, k5x5 & $16 \times 16$    & 112  & 3 \\
6  & MBConv6, k5x5 & $16 \times 16$    & 192  & 4 \\
7  & MBConv6, k3x3 & $8 \times 8$      & 320  & 1 \\
8  & Conv1x1       & $8 \times 8$      & 1280 & 1 \\
\bottomrule
\end{tabular}
\end{table}

\section{Flujo de tensores}

\begin{minted}[language=Python, basicstyle=\ttfamily\small, frame=single]{python}
# Entrada: imagen preprocesada (normalizada)
x:  (B, 3, 512, 512)

# EfficientNet-B0
x = visual_network(x)
# -> (B, 1280, 16, 16)  <- etapa 9, depende del tamaño de entrada

# DSC 1280 -> 1024
x = visual_features_downscale(x)
# -> (B, 1024, 16, 16)

# AdaptiveMaxPool2d((1, 1))
x = adaptive_max_pooling(x)
# -> (B, 1024, 1, 1)

# DSC final 1024 -> 1024
x = depthwise_seperable_convolution(x)
# -> (B, 1024, 1, 1)

# squeeze última dim
x = x.squeeze(-1).contiguous()
# -> (B, 1024, 1)
\end{minted}

\section{Parámetros}

\begin{table}[H]
\centering
\caption{Parámetros por componente visual.}
\label{tab:vis-params-linea}
\begin{tabular}{lr}
\toprule
\textbf{Componente} & \textbf{Parámetros} \\
\midrule
\texttt{EfficientNet-B0 features}   & \textasciitilde 5,29 M \\
\texttt{DSC(1280, 1024)}             & \textasciitilde 14 K \\
\texttt{AdaptiveMaxPool2d(1,1)}      & 0 \\
\texttt{DSC(1024, 1024)}             & \textasciitilde 10 K \\
\midrule
\textbf{Total visual}                & \textbf{\textasciitilde 5,31 M} \\
\bottomrule
\end{tabular}
\end{table}

% ============================================================
% Capítulo 7 — Language Network: conceptos + código
% ============================================================
\chapter{Red de lenguaje: BERT, embeddings, BiGRU y pack-padded}
\label{cap:language-network}

Este capítulo explica los conceptos de NLP detrás de la rama
textual del modelo y luego muestra el código línea por línea.

\section{¿Qué es BERT y por qué sólo su tokenizer?}

\subsection{BERT}

BERT (\emph{Bidirectional Encoder Representations from
Transformers})~\cite{devlin2018bert} es un modelo de lenguaje
preentrenado en un corpus enorme de texto. Tiene dos partes:

\begin{enumerate}
  \item \textbf{Tokenizer}: convierte texto a IDs numéricos usando
        un vocabulario fijo. Es lo único que usamos aquí.
  \item \textbf{Transformer encoder}: 12 capas de atención que
        producen embeddings contextuales. \textbf{No lo usamos}.
\end{enumerate}

\subsection{¿Por qué no usar BERT entero?}

BERT completo tiene \textasciitilde 110M parámetros. BiGRU sólo
\textasciitilde 4.7M. Para textos cortos (descripciones de 30
palabras) BiGRU captura suficiente contexto. Esto reduce:

\begin{itemize}
  \item Coste computacional (\textasciitilde 20$\times$ menos).
  \item VRAM (cabe en GPU modestas).
  \item Tiempo de entrenamiento.
\end{itemize}

\subsection{Cómo funciona el tokenizer de BERT}

BERT usa \emph{WordPiece tokenization}: el vocabulario está
formado por sub-palabras (\emph{tokens}) frecuentes más letras
individuales. Ejemplo:

\begin{verbatim}
Entrada:  "the man is wearing a black jacket"
Tokens:   ['the', 'man', 'is', 'wearing', 'a', 'black', 'jacket']
IDs:      [1996, 2158, 2003, 7078, 1037, 2304, 7505]
Especial: [CLS] = 101, [SEP] = 102
\end{verbatim}

\noindent\textbf{Vocabulario}: 30\,522 tokens base. Más 1 (índice
0 reservado para padding) $\approx$ 30\,609 únicos. Más 1 para
\texttt{[PAD]} = 29\,610 (lo que usa este proyecto).

\section{¿Qué es un embedding?}

Un \emph{embedding} es una representación densa (vector de
números reales) de un token. En este proyecto:

\begin{itemize}
  \item Cada token del vocabulario $\to$ vector de 512 números
        reales.
  \item Total de parámetros: $29\,610 \times 512 = 15,16$M.
\end{itemize}

\noindent\textbf{Inicialización}: $\mathcal{N}(0, 1)$ por defecto.
Durante el entrenamiento, estos vectores se ajustan para que
tokens con significado similar tengan vectores cercanos.

\section{¿Qué es un padding?}

Los textos del batch tienen longitudes distintas. Para meterlos
en un tensor uniforme se \emph{padean} a una longitud fija (100
en este proyecto):

\begin{verbatim}
"a man in red"     -> [101, 1037, 2158, 1999, 2417, 102] (6 tokens)
"a tall person"    -> [101, 1037, 4206, 2718, 102] (5 tokens)
\end{verbatim}

\noindent Ambos se convierten a longitud 100 añadiendo ceros
(\texttt{[PAD]} = 0) al final. El modelo \textbf{ignora} el padding
porque:

\begin{enumerate}
  \item El índice 0 nunca se usa como token real (es
        \texttt{[PAD]}).
  \item El embedding de índice 0 está fijo en $(0, 0, \ldots, 0)$.
  \item La GRU recibe \texttt{orig\_token\_length} para saber
        dónde termina cada secuencia.
\end{enumerate}

\section{¿Qué es una GRU?}

Una \emph{Gated Recurrent Unit}~\cite{cho2014gru} es una versión
simplificada de LSTM. Procesa una secuencia token por token y
mantiene un \emph{estado oculto} $h_t$:

\begin{align}
  z_t &= \sigma(W_z \cdot [h_{t-1}, x_t])  &\text{(update gate)} \\
  r_t &= \sigma(W_r \cdot [h_{t-1}, x_t])  &\text{(reset gate)} \\
  \tilde{h}_t &= \tanh(W \cdot [r_t \odot h_{t-1}, x_t]) \\
  h_t &= (1 - z_t) \odot h_{t-1} + z_t \odot \tilde{h}_t
\end{align}

\noindent\textbf{Bidireccional} significa que se corren dos GRU en
paralelo: una hacia adelante y otra hacia atrás. Las salidas se
\emph{promedian} para tener en cuenta contexto en ambas
direcciones.

\noindent\textbf{Ventajas} sobre LSTM:
\begin{itemize}
  \item Menos parámetros (2 compuertas vs 3).
  \item Más rápido en secuencias cortas.
\end{itemize}

\section{¿Qué es pack-padded-sequence?}

\subsection{El problema}

Si procesamos un batch con secuencias de longitudes [10, 30, 5,
20] todas padeadas a 100, la GRU hace trabajo inútil en los
\texttt{[PAD]}. Además los \texttt{[PAD]} afectan el estado
oculto: como $h_{100} = f(h_{99}, 0)$, el estado final está
"contaminado" por 70 pasos de padding.

\subsection{La solución}

\texttt{pack\_padded\_sequence} de PyTorch:

\begin{enumerate}
  \item \textbf{Ordena} las secuencias por longitud descendente.
  \item \textbf{Elimina} los pasos de padding antes de pasar a
        la RNN.
  \item La RNN procesa sólo los tokens reales.
  \item \texttt{pad\_packed\_sequence} restaura el orden original
        con padding.
\end{enumerate}

\noindent\textbf{Requisito}: pasar las longitudes reales en CPU
(no GPU) porque el backend de PyTorch así lo requiere.

\section{Código: \texttt{GRULanguageNetwork}}

\begin{minted}[language=Python, numbers, firstnumber=11, linerange=11-41]{python}
class GRULanguageNetwork(nn.Module):
    """Embedding + Dropout + GRU bidireccional."""

    def __init__(self, config=None):
        super(GRULanguageNetwork, self).__init__()
        if config is None:
            raise ValueError("`config` can not be none")

        # Capa de embedding
        # vocab_size = 29609 (BERT) + 1 (padding) = 29610
        # embedding_dim = 512
        # padding_idx = 0: el vector para índice 0 se mantiene en 0
        self.embedding = nn.Embedding(
            config.vocab_size,        # 29610
            config.embedding_dim,     # 512
            padding_idx=0
        )

        # Dropout tras el embedding (30%)
        self.dropout = nn.Dropout(config.dropout_rate)  # 0.30

        # GRU bidireccional
        # input_size: 512 (dim del embedding)
        # hidden_size: 1024 (feature_length)
        # num_layers: 1
        # bidirectional=True: procesa en ambas direcciones
        # bias=False: ahorra parámetros
        self.gru = nn.GRU(
            config.embedding_dim,     # 512
            config.feature_length,    # 1024
            num_layers=config.num_layers,    # 1
            bidirectional=True,
            bias=False
        )

    def forward(self, text_ids=None, text_length=None):
        """text_ids: (B, 100), text_length: (B,).
           returns: (B, 1024, 100, 1)"""
        # 1. Embedding: (B, 100) -> (B, 100, 512)
        text_embedding = self.embedding(text_ids)
        # 2. Dropout (solo en training mode)
        text_embedding = self.dropout(text_embedding)
        # 3. Procesar con BiGRU (usando pack_padded_sequence)
        feature = self.do_RNN_computation(text_embedding, text_length, self.gru)
        return feature
\end{minted}

\section{Código: \texttt{do\_RNN\_computation}}

\begin{minted}[language=Python, numbers, firstnumber=44, linerange=44-69]{python}
    def do_RNN_computation(self, text_embedding, text_length, gru):
        """Procesa la secuencia con GRU usando pack_padded_sequence."""

        # Asegurar que text_length es 1-D
        text_length = text_length.view(-1)

        # Ordenar índices por longitud DESCENDENTE
        _, sort_index = torch.sort(text_length, dim=0, descending=True)

        # Invertir el orden: nos servirá para "des-ordenar" después
        _, unsort_index = sort_index.sort()

        # Reordenar el batch por longitud descendente
        sortlength_text_embedding = text_embedding[sort_index, :]
        sort_text_length = text_length[sort_index]

        # Empaquetar: ignora el padding en el cómputo de la GRU
        packed_text_embedding = pack_padded_sequence(
            sortlength_text_embedding,
            sort_text_length.cpu(),
            batch_first=True
        )

        # Pasar por la GRU
        packed_feature, hn = gru(packed_text_embedding)

        # Desempaquetar: vuelve a la forma (B, T, *)
        total_length = text_embedding.size(1)
        sort_feature = pad_packed_sequence(
            packed_feature, batch_first=True, total_length=total_length
        )

        # Des-ordenar: volver al orden original del batch
        unsort_feature = sort_feature[0][unsort_index, :]

        # Promediar fwd y bwd:
        # BiGRU output: (B, T, 2*hidden) = (B, 100, 2048)
        # [:, :, :1024] = forward, [:, :, 1024:] = backward
        unsort_feature = (
            unsort_feature[:, :, :int(unsort_feature.size(2) / 2)] +
            unsort_feature[:, :, int(unsort_feature.size(2) / 2):]
        ) / 2

        # Preparar para alimentar el DSC final:
        # permute(0,2,1): (B, 100, 1024) -> (B, 1024, 100)
        # unsqueeze(3): añadir dim al final -> (B, 1024, 100, 1)
        return unsort_feature.permute(0, 2, 1).contiguous().unsqueeze(3)
\end{minted}

\section{Código: \texttt{BERTTokenizer}}

\begin{minted}[language=Python, numbers, firstnumber=10, linerange=10-19]{python}
class BERTTokenizer(object):
    """Adapta el tokenizer de HuggingFace a la interfaz del proyecto."""

    def __init__(self):
        super(BERTTokenizer, self).__init__()

        # Importar transformers
        _, tokenizer_class, pretrained_weights = (
            ppb.BertModel, ppb.BertTokenizer, 'bert-base-uncased'
        )

        # Cargar el tokenizer preentrenado (BERT-base, vocab inglés)
        self.tokenizer = tokenizer_class.from_pretrained(pretrained_weights)

    def __call__(self, text: str = None) -> torch.LongTensor:
        """Codifica texto a IDs de tokens."""
        # .encode() añade automáticamente [CLS] y [SEP]
        tokens = self.tokenizer.encode(text)
        return tokens
\end{minted}

\section{Código: \texttt{pad\_tokens}}

\begin{minted}[language=Python, numbers, firstnumber=86, linerange=86-107]{python}
def pad_tokens(tokens, tokens_length_max):
    """Padea o trunca tokens a longitud fija."""
    tokens_length = len(tokens)

    # Convertir a tensor 1-D
    tokens = torch.from_numpy(np.array(tokens)).view(-1).contiguous().float()

    if tokens_length < tokens_length_max:
        # Caso corto: añadir zeros de padding al final
        zero_padding = torch.zeros(tokens_length_max - tokens_length)
        tokens = torch.cat([tokens, zero_padding], 0)
    else:
        # Caso largo: truncar a tokens_length_max
        tokens = tokens[:tokens_length_max]
        tokens_length = tokens_length_max

    return tokens, tokens_length
\end{minted}

\section{Flujo de tensores}

\begin{minted}[language=Python, basicstyle=\ttfamily\small, frame=single]{python}
# Entrada
text_ids:     (B, 100)    LongTensor, padded con 0s
text_length:  (B,)        longitudes reales (5-50 típico)

# 1. Embedding
out = embedding(text_ids)
# (B, 100) -> (B, 100, 512)

# 2. Dropout (solo training)
out = dropout(out)
# (B, 100, 512)

# 3. BiGRU con pack_padded_sequence
out = do_RNN_computation(out, text_length, self.gru)
# (B, 100, 512) -> (B, 100, 2048) -> (B, 100, 1024) -> (B, 1024, 100, 1)

# 4. Max-pool sobre tokens (en text_embedding)
out = max(out, dim=2)
# (B, 1024, 1, 1)

# 5. DSC compartido
out = depthwise_seperable_convolution(out)
# (B, 1024, 1, 1)

# 6. squeeze
out = out.squeeze(-1).contiguous()
# (B, 1024, 1)
\end{minted}

\section{Parámetros}

\begin{table}[H]
\centering
\caption{Parámetros de la red de lenguaje.}
\label{tab:lang-params-linea}
\begin{tabular}{lr}
\toprule
\textbf{Componente} & \textbf{Parámetros} \\
\midrule
\texttt{Embedding(29610, 512)} & 15,16 M \\
\texttt{Dropout(0.30)}         & 0 \\
\texttt{BiGRU(512→1024, bidir)} & 4,72 M \\
\texttt{DSC(1024→1024)}        & 10 K \\
\midrule
\textbf{Total lenguaje}        & \textbf{19,89 M} \\
\bottomrule
\end{tabular}
\end{table}

% ============================================================
% Capítulo 8 — TextReIDNet: conceptos + código
% ============================================================
\chapter{TextReIDNet: joint embedding visual-textual}
\label{cap:textreidnet}

Este capítulo explica los conceptos detrás del modelo principal y
muestra su código línea por línea.

\section{¿Qué es un joint embedding?}

\subsection{El concepto}

Un \emph{joint embedding space} (espacio de embedding conjunto)
es un espacio vectorial donde las dos modalidades --- imágenes
y texto --- se mapean al mismo lugar. La idea clave:

\begin{itemize}
  \item Imagen de una persona $\to$ vector $\mathbf{I}' \in
        \mathbb{R}^{1024}$.
  \item Descripción de esa persona $\to$ vector $\mathbf{T}' \in
        \mathbb{R}^{1024}$.
  \item \textbf{Si el modelo es bueno, los vectores están
        \emph{cerca} en el espacio}.
\end{itemize}

\noindent\textbf{Entonces} medir qué tan "cerca" están dos
vectores nos dice qué tan parecidas son las dos entradas.

\section{¿Qué es la similitud coseno?}

Dados dos vectores $\mathbf{a}, \mathbf{b} \in \mathbb{R}^n$:

\begin{equation}
  \text{sim}(\mathbf{a}, \mathbf{b}) = \frac{\mathbf{a} \cdot \mathbf{b}}
                                            {\|\mathbf{a}\| \cdot \|\mathbf{b}\|}
  \in [-1, +1]
  \label{eq:cosine}
\end{equation}

\noindent\textbf{Interpretación}:
\begin{itemize}
  \item $+1$: vectores idénticos (apuntan en la misma dirección).
  \item $0$: ortogonales (no relacionados).
  \item $-1$: opuestos.
\end{itemize}

\noindent\textbf{Por qué coseno y no distancia euclidiana?}:
\begin{itemize}
  \item Invariante a la magnitud: $\mathbf{a}$ y $2\mathbf{a}$
        tienen similitud 1.
  \item Es un \emph{ángulo} entre vectores, no una distancia.
  \item Valores acotados en $[-1, +1]$, fácil de interpretar.
\end{itemize}

\section{¿Por qué un DSC compartido?}

La última capa (\texttt{DepthwiseSeparableConv(1024, 1024)}) se
usa para \textbf{ambas} ramas (visual y textual). Esto tiene dos
efectos:

\begin{enumerate}
  \item \textbf{Menos parámetros}: una capa en lugar de dos.
  \item \textbf{Mismo espacio}: las dos modalidades se proyectan
        con la misma transformación, lo que \emph{fuerza} a que
        vivan en el mismo espacio.
\end{enumerate}

\noindent\textbf{Analogía}: es como si dos personas midieran en
metros y otra persona midiera en pies. Si no comparten unidad,
no se pueden comparar. El DSC compartido es la "conversión a la
misma unidad".

\section{¿Qué hace \texttt{squeeze} y \texttt{contiguous}?}

\subsection{\texttt{squeeze(dim)}}

Elimina dimensiones de tamaño 1. Por ejemplo:
\begin{itemize}
  \item $(B, 1024, 1, 1) \xrightarrow{\text{squeeze}(-1)} (B,
        1024, 1)$.
  \item $(B, 1024, 1) \xrightarrow{\text{squeeze}(-1)} (B,
        1024)$.
\end{itemize}

\noindent El argumento $-1$ indica "la última dimensión". Útil
para preparar tensores para operaciones que esperan un número
específico de dims.

\subsection{\texttt{contiguous()}}

Algunos operaciones de PyTorch (especialmente \texttt{view},
\texttt{reshape}, slicing) requieren que el tensor sea
\emph{contiguo} en memoria. Si no lo es, se hace una copia
interna. Llamar \texttt{.contiguous()} fuerza la contigüidad
explícitamente.

\section{¿Por qué dimensión 1024?}

La dimensión del espacio conjunto ($1024$) es un compromiso:

\begin{itemize}
  \item \textbf{Suficientemente grande} para capturar la
        semántica de ambas modalidades.
  \item \textbf{Suficientemente pequeño} para:
        \begin{itemize}
          \item Similitud coseno estable numéricamente.
          \item Inferencia rápida (4 KB por embedding en FP32).
          \item Cabe en GPUs modestas.
        \end{itemize}
  \item \textbf{Coincide} con las unidades ocultas de la BiGRU,
        lo que permite compartir el último DSC sin mismatch de
        dimensiones.
\end{itemize}

\section{Código: Constructor \texttt{\_\_init\_\_}}

\begin{minted}[language=Python, numbers, firstnumber=13, linerange=13-44]{python}
class TextReIDNet(nn.Module):
    """La clase principal del modelo."""

    def __init__(self, configs: dict = None):
        # Llama al constructor de nn.Module (registra parámetros,
        # buffers, submódulos, etc.)
        super(TextReIDNet, self).__init__()

        # Validación: configs es obligatorio
        if configs is None:
            raise ValueError("`configs` parameter can not be None")

        # Guardar referencia a la configuración completa
        self.configs = configs

        # --- Rama visual ---
        # VisualNetwork = EfficientNet-B0 + wrapper
        self.visual_network = VisualNetwork()

        # --- Rama de lenguaje ---
        # GRULanguageNetwork = Embedding + Dropout + BiGRU
        self.language_network = GRULanguageNetwork(configs)

        # --- Capas de proyección al espacio conjunto ---
        # 1ra DSC: reduce canales de 1280 (EffNet) a 1024 (joint dim)
        self.visual_features_downscale = DepthwiseSeparableConv(1280, 1024)

        # Adaptive max pooling: colapsa (H, W) a (1, 1)
        self.adaptive_max_pooling = nn.AdaptiveMaxPool2d((1, 1))

        # 2da DSC: mantiene 1024 canales, proyecta al espacio conjunto
        # COMPARTIDA con la rama textual (ver text_embedding)
        self.depthwise_seperable_convolution = DepthwiseSeparableConv(
            1024, self.configs.feature_length
        )
\end{minted}

\section{Código: \texttt{forward}}

\begin{minted}[language=Python, numbers, firstnumber=45, linerange=45-59]{python}
    def forward(self, image, text_ids=None, text_length=None):
        """image: (B,3,H,W), text_ids: (B,100), text_length: (B,)
           returns: ((B,1024,1), (B,1024,1))"""
        # Calcular embedding visual (rama izquierda)
        visaual_embeddings = self.image_embedding(image)

        # Calcular embedding textual (rama derecha)
        textual_embeddings = self.text_embedding(text_ids, text_length)

        # Devolver ambos. train.py los aplana a (B, 1024) con view(-1).
        # Nota: hay un typo en el código: "visaual" en lugar de "visual"
        return visaual_embeddings, textual_embeddings
\end{minted}

\section{Código: \texttt{image\_embedding}}

\begin{minted}[language=Python, numbers, firstnumber=62, linerange=62-72]{python}
    def image_embedding(self, image=None):
        """image: (B, 3, H, W) batch de imágenes.
           returns: (B, 1024, 1) embeddings visuales."""

        # PASO 1: EfficientNet-B0
        # Entrada: (B, 3, 512, 512)
        # Salida: (B, 1280, 16, 16) <- última etapa
        visual_features = self.visual_network(image)

        # PASO 2: DSC 1280 -> 1024
        # Entrada: (B, 1280, 16, 16)
        # Salida:  (B, 1024, 16, 16)
        visual_features = self.visual_features_downscale(visual_features)

        # PASO 3: Adaptive Max Pooling
        # Entrada: (B, 1024, 16, 16)
        # Salida:  (B, 1024, 1, 1)
        visual_features = self.adaptive_max_pooling(visual_features)

        # PASO 4: DSC final 1024 -> 1024 (compartido con texto)
        # Entrada: (B, 1024, 1, 1)
        # Salida:  (B, 1024, 1, 1)
        visual_features = self.depthwise_seperable_convolution(visual_features)

        # PASO 5: squeeze + contiguous
        # .squeeze(-1) elimina la última dim (1) -> (B, 1024, 1)
        # .contiguous() asegura memoria contigua
        visual_features = visual_features.squeeze(-1).contiguous()
        return visual_features
\end{minted}

\section{Código: \texttt{text\_embedding}}

\begin{minted}[language=Python, numbers, firstnumber=75, linerange=75-89]{python}
    def text_embedding(self, text_ids=None, text_length=None):
        """text_ids: (B, 100), text_length: (B,).
           returns: (B, 1024, 1) embeddings textuales."""

        # PASO 1: Pasar por la red de lenguaje
        # Salida de GRULanguageNetwork:
        # Entrada:  (B, 100)         tokens BERT
        #   Embedding: (B, 100, 512)
        #   BiGRU:     (B, 100, 2048) bidireccional
        #   avg f+b:   (B, 100, 1024) promedio de fwd + bwd
        #   permute:   (B, 1024, 100, 1)
        textual_features = self.language_network(text_ids, text_length)

        # PASO 2: Max-pooling sobre la dim de tokens (dim=2)
        # Para cada canal (1024), tomar el valor máximo sobre los 100 tokens
        # Entrada: (B, 1024, 100, 1)
        # Salida:  (B, 1024, 1, 1)
        textual_features, _ = torch.max(textual_features, dim=2, keepdim=True)

        # PASO 3: DSC compartido con la rama visual
        textual_features = self.depthwise_seperable_convolution(textual_features)

        # PASO 4: squeeze + contiguous
        textual_features = textual_features.squeeze(-1).contiguous()
        return textual_features
\end{minted}

\section{Diagrama del flujo completo}

\begin{figure}[H]
\centering
\begin{tikzpicture}[
    font=\sffamily\footnotesize,
    node distance=0.7cm,
    box/.style={
        rectangle, draw=black, thick, rounded corners=2pt,
        minimum height=0.8cm, minimum width=2.2cm,
        align=center, fill=blue!10
    },
    arr/.style={-{Stealth[length=2.5mm,width=2mm]}, thick}
]
\node[box] (i) {Imagen\\$(B, 3, 512, 512)$};
\node[box, below=of i] (eff) {EffNet-B0\\$(B, 1280, 16, 16)$};
\node[box, below=of eff] (d1) {DSC\\$1280 \to 1024$};
\node[box, below=of d1] (p) {Pool $(1,1)$};
\node[box, below=of p] (d2) {DSC\\$1024 \to 1024$};
\node[box, below=of d2, fill=green!15] (ve) {$\mathbf{I}' (B, 1024, 1)$};

\node[box, right=3cm of i] (t) {Tokens\\$(B, 100)$};
\node[box, right=3cm of eff] (e) {Embedding\\$(B, 100, 512)$};
\node[box, right=3cm of d1] (g) {BiGRU\\$(B, 100, 2048)$};
\node[box, right=3cm of d2] (m) {Max$+$Perm\\$(B, 1024, 1, 1)$};

\draw[arr] (i) -- (eff);
\draw[arr] (eff) -- (d1);
\draw[arr] (d1) -- (p);
\draw[arr] (p) -- (d2);
\draw[arr] (d2) -- (ve);

\draw[arr] (t) -- (e);
\draw[arr] (e) -- (g);
\draw[arr] (g) -- (m);
\draw[arr] (m.east) -| ([yshift=0.4cm]ve.east) -- (ve.east);

\draw[arr, dashed] (ve) -- node[right] {cosine} +(2, 0);
\end{tikzpicture}
\caption{Flujo completo de tensores. Notar que el último DSC es
compartido entre ambas ramas.}
\label{fig:textreidnet-flow}
\end{figure}

\section{Inicialización de pesos}

\texttt{TextReIDNet} no llama ninguna función de inicialización.
Los pesos heredan:

\begin{itemize}
  \item \textbf{EfficientNet-B0}: preentrenado en ImageNet.
  \item \textbf{Embedding}: $\mathcal{N}(0, 1)$ por defecto de
        PyTorch.
  \item \textbf{BiGRU}: inicialización por defecto de PyTorch.
  \item \textbf{DSC}: Kaiming (\texttt{weights\_init\_kaiming}).
\end{itemize}

\section{Inferencia: extraer embeddings}

\begin{minted}[language=Python, basicstyle=\ttfamily\small, frame=single]{python}
import torch
from config import sys_configuration
from model.textreidnet import TextReIDNet

cfg = sys_configuration(dataset_source='huggingface')
model = TextReIDNet(cfg).to(cfg.device).eval()

# Modo evaluación (desactiva dropout)
with torch.no_grad():
    # Embedding de imagen
    img = torch.randn(2, 3, 512, 512).to(cfg.device)
    img_emb = model.image_embedding(img)        # (2, 1024, 1)
    img_emb = img_emb.view(img_emb.size(0), -1)  # (2, 1024)

    # Embedding de texto
    txt = torch.zeros(2, 100, dtype=torch.long).to(cfg.device)
    lng = torch.tensor([50, 30]).to(cfg.device)
    txt_emb = model.text_embedding(txt, lng)     # (2, 1024, 1)
    txt_emb = txt_emb.view(txt_emb.size(0), -1)  # (2, 1024)

    # Similitud coseno entre las 2 imágenes y los 2 textos
    img_emb_n = img_emb / img_emb.norm(dim=1, keepdim=True)
    txt_emb_n = txt_emb / txt_emb.norm(dim=1, keepdim=True)
    sim = torch.mm(img_emb_n, txt_emb_n.t())     # (2, 2)
    print(sim)
\end{minted}

% ============================================================
% Capítulo 9 — Losses: conceptos + código
% ============================================================
\chapter{Funciones de pérdida: triplet y cross-entropy}
\label{cap:losses}

Este capítulo explica los conceptos detrás de las dos pérdidas
usadas y muestra su código línea por línea.

\section{¿Qué es un embedding y por qué necesita pérdidas?}

Después del forward tenemos dos vectores por muestra:
\begin{itemize}
  \item $\mathbf{I}' \in \mathbb{R}^{1024}$ (imagen).
  \item $\mathbf{T}' \in \mathbb{R}^{1024}$ (texto).
\end{itemize}

\noindent\textbf{Objetivo}: vectores de la misma persona deben
estar \emph{cerca}, y de personas diferentes \emph{lejos}.

\noindent\textbf{Problema}: el modelo no sabe cómo hacerlo solo.
Necesitamos una función que mida qué tan mal lo está haciendo y
guíe el aprendizaje: la \textbf{pérdida} (loss).

\section{¿Qué es la triplet loss?}

\subsection{Concepto}

La \emph{triplet loss}~\cite{schroff2015facenet} trabaja con
\emph{triplets}: tríos de muestras (anchor, positive, negative).

\begin{itemize}
  \item \textbf{Anchor} (A): muestra de referencia.
  \item \textbf{Positive} (P): muestra del \emph{mismo} pid.
  \item \textbf{Negative} (N): muestra de \emph{otro} pid.
\end{itemize}

La pérdida enforces que:

\begin{equation}
  \text{sim}(A, P) - \text{sim}(A, N) + \text{margin} \leq 0
  \label{eq:triplet}
\end{equation}

\noindent\textbf{En palabras}: el anchor debe estar \emph{más
cerca} del positive que del negative, con un margen mínimo
($\text{margin} = 0.5$ en este proyecto).

\subsection{¿Por qué el margen?}

Sin margen, basta con que $\text{sim}(A, P) > \text{sim}(A, N)$.
Esto da una solución trivial: todos los embeddings en un punto.
Con margen exigimos que la diferencia sea $\geq \text{margin}$,
lo que \emph{obliga} a aprender embeddings informativos.

\section{¿Qué es semi-hard mining?}

\subsection{Los tres tipos de negativos}

Para un anchor dado, hay tres tipos de negativos:

\begin{enumerate}
  \item \textbf{Easy}: $\text{sim}(A, N) \ll \text{sim}(A, P)$.
        El modelo ya los clasifica bien. \emph{No ayudan} al
        aprendizaje.
  \item \textbf{Hard}: $\text{sim}(A, N) > \text{sim}(A, P)$.
        El modelo ya los confunde. Contribuyen mucho al
        gradiente pero pueden hacer el entrenamiento inestable.
  \item \textbf{Semi-hard}: $\text{sim}(A, P) < \text{sim}(A, N)
        < \text{sim}(A, P) + \text{margin}$. El modelo está
        cerca pero no del todo mal. \emph{Ideal}: gradiente
        útil + entrenamiento estable.
\end{enumerate}

\subsection{Criterio usado en este proyecto}

El código implementa semi-hard mining~\cite{schroff2015facenet}:
\begin{itemize}
  \item Si hay negativos semi-hard disponibles, elegir uno
    aleatoriamente.
  \item Si no hay ninguno, usar uno cualquiera (caso
    degenerado).
\end{itemize}

\noindent\textbf{Ventaja}: estable y eficiente; converge más
rápido que easy mining sin los problemas de hard mining.

\section{¿Qué es la cross-entropy?}

\subsection{Clasificación como identidad}

Otra forma de entrenar el modelo es agregar una cabeza de
\textbf{clasificación} que predice el pid (identidad de persona).
Si el modelo predice correctamente, está discriminando bien.

La cross-entropy mide qué tan lejos están las predicciones del
label correcto:

\begin{equation}
  \mathcal{L}_{\text{CE}} = -\sum_{c=1}^{C} y_c \cdot \log(\hat{y}_c)
  \label{eq:ce}
\end{equation}

\noindent donde $y_c$ es 1 si $c$ es el pid correcto, $\hat{y}_c$ es
la probabilidad predicha para la clase $c$.

\subsection{PyTorch lo hace solo}

\texttt{nn.CrossEntropyLoss(reduction='mean')} computa la CE
promediada sobre el batch. Internamente aplica
\texttt{LogSoftmax + NLLLoss}.

\section{¿Por qué simétrica (i2t + t2i)?}

Tanto RankingLoss como IdentityLoss se calculan en
\textbf{ambas direcciones}:
\begin{itemize}
  \item \textbf{i2t}: predice la clase del texto desde el
    embedding de la imagen.
  \item \textbf{t2i}: predice la clase de la imagen desde el
    embedding del texto.
\end{itemize}

\noindent\textbf{¿Por qué?} El modelo debe ser bueno
\emph{simétricamente}: si una imagen $\leftrightarrow$ texto
describe a la misma persona, ambos sentidos deben predecir el
mismo pid. Esto fuerza coherencia en ambos lados del embedding.

\section{Pérdida total}

\begin{equation}
  \mathcal{L}_{\text{total}} = \alpha \cdot \mathcal{L}_{\text{rank}}
                             + \beta  \cdot \mathcal{L}_{\text{ID}}
\end{equation}

\noindent En este proyecto $\alpha = \beta = 1.0$
(\fline{config.py}{87-88}).

\section{Código: \texttt{calculate\_similarity}}

\begin{minted}[language=Python, numbers, firstnumber=11, linerange=11-18]{python}
def calculate_similarity(image_embedding, text_embedding):
    """Calcula la matriz de similitud coseno entre imágenes y textos."""

    # Aplanar embeddings a (B, 1024) por si vienen con más dims
    image_embedding = image_embedding.view(image_embedding.size(0), -1)
    text_embedding = text_embedding.view(text_embedding.size(0), -1)

    # Normalizar L2 (norma = 1)
    # +1e-8 para evitar división por cero
    image_embedding_norm = image_embedding / (
        image_embedding.norm(dim=1, keepdim=True) + 1e-8
    )
    text_embedding_norm = text_embedding / (
        text_embedding.norm(dim=1, keepdim=True) + 1e-8
    )

    # Producto matricial: similitud entre cada par (img_i, txt_j)
    similarity = torch.mm(image_embedding_norm, text_embedding_norm.t())
    return similarity
\end{minted}

\section{Código: \texttt{RankingLoss}}

\begin{minted}[language=Python, numbers, firstnumber=21, linerange=21-51]{python}
class RankingLoss(nn.Module):
    """Triplet loss con minería semi-hard."""

    def __init__(self, config=None):
        super(RankingLoss, self).__init__()
        self.config = config

    def semi_hard_negative(self, loss):
        """Elige un negativo semi-hard al azar."""
        # loss: array con losses para cada negativo
        # logical_and(0 < loss < margin)
        negative_index = np.where(
            np.logical_and(loss < self.config.margin, loss > 0)
        )[0]
        return np.random.choice(negative_index) if len(negative_index) > 0 else None

    def get_triplets(self, similarity, labels):
        """Para cada anchor, busca un par (positivo, negativo)."""
        similarity = similarity.cpu().data.numpy()
        labels = labels.cpu().data.numpy()
        triplets = []

        for idx, label in enumerate(labels):
            # Índices con OTRO pid (negativos)
            negative = np.where(labels != label)[0]
            # Similitud del anchor consigo mismo (~1)
            ap_sim = similarity[idx, idx]
            # loss = sim(a,n) - sim(a,p) + margin
            loss = similarity[idx, negative] - ap_sim + self.config.margin
            # Elegir un negativo semi-hard
            negetive_index = self.semi_hard_negative(loss)
            if negetive_index is not None:
                triplets.append([idx, idx, negative[negetive_index]])

        if len(triplets) == 0:
            triplets.append([idx, idx, negative[0]])
        triplets = np.array(triplets)
        return torch.LongTensor(triplets)
\end{minted}

\begin{minted}[language=Python, numbers, firstnumber=54, linerange=54-72]{python}
    def forward(self, img, txt, label, sim_neg=None):
        """img: (B,1024), txt: (B,1024), label: (B,)"""
        similarity = calculate_similarity(img, txt)
        if sim_neg is not None:
            similarity = similarity + sim_neg

        # Triplets image->text
        image_triplets = self.get_triplets(similarity, label)
        image_anchor_loss = F.relu(
            self.config.margin
            - similarity[image_triplets[:, 0], image_triplets[:, 1]]
            + similarity[image_triplets[:, 0], image_triplets[:, 2]]
        )

        # Triplets text->image (transponer)
        text_triplets = self.get_triplets(similarity.t(), label)
        similarity = similarity.t()
        texy_anchor_loss = F.relu(
            self.config.margin
            - similarity[text_triplets[:, 0], text_triplets[:, 1]]
            + similarity[text_triplets[:, 0], text_triplets[:, 2]]
        )

        loss = torch.sum(image_anchor_loss) + torch.sum(texy_anchor_loss)
        return loss
\end{minted}

\section{Código: \texttt{IdentityLoss}}

\begin{minted}[language=Python, numbers, firstnumber=10, linerange=10-27]{python}
def weights_init_classifier(m):
    """Inicializa capas Linear con N(0, 0.001) y bias=0."""
    classname = m.__class__.__name__
    if classname.find('Linear') != -1:
        init.normal_(m.weight.data, std=0.001)
        init.constant_(m.bias.data, 0.0)


class classifier(nn.Module):
    """Capa lineal simple: input_dim -> output_dim."""

    def __init__(self, input_dim, output_dim):
        super(classifier, self).__init__()
        self.block = nn.Linear(input_dim, output_dim)
        self.block.apply(weights_init_classifier)

    def forward(self, x):
        return self.block(x)
\end{minted}

\begin{minted}[language=Python, numbers, firstnumber=30, linerange=30-56]{python}
class IdentityLoss(nn.Module):
    """Cross-entropy simétrica: image->text + text->image."""

    def __init__(self, config=None, part=1, class_num=None):
        super(IdentityLoss, self).__init__()
        self.part = part

        # Crear `part` clasificadores (en este proyecto part=1)
        W = []
        for i in range(part):
            W.append(classifier(config.feature_length, class_num))
        self.W = nn.Sequential(*W)

    def calculate_IdLoss(self, image_embedding_local,
                         text_embedding_local, label):
        """image_embedding_local: (B,1024,1), text_embedding_local: (B,1024,1)"""
        label = label.view(label.size(0))
        criterion = nn.CrossEntropyLoss(reduction='mean')

        Lipt_local = 0  # image->text loss
        Ltpi_local = 0  # text->image loss

        for i in range(self.part):
            # Proyectar embedding al espacio de clases
            score_i2t_local_i = self.W[i](image_embedding_local[:, :, i])
            score_t2i_local_i = self.W[i](text_embedding_local[:, :, i])

            Lipt_local += criterion(score_i2t_local_i, label)
            Ltpi_local += criterion(score_t2i_local_i, label)

        loss = (Lipt_local + Ltpi_local) / self.part
        return loss

    def forward(self, image_embedding_local,
                text_embedding_local, label):
        return self.calculate_IdLoss(image_embedding_local,
                                     text_embedding_local, label)
\end{minted}

\section{Uso conjunto en \texttt{train.py}}

\begin{minted}[language=Python, numbers, firstnumber=130, linerange=130-150]{python}
with torch.autocast(device_type=config.device, dtype=torch.float16):
    # Forward
    visual_embeddings, textual_embeddings = model(
        image=preprocessed_images,
        text_ids=token_ids,
        text_length=orig_token_length
    )

    # Aplanar a (B, 1024) - RankingLoss espera aplanados
    visual_emb_flat = visual_embeddings.view(visual_embeddings.size(0), -1)
    text_emb_flat = textual_embeddings.view(textual_embeddings.size(0), -1)
    labels_flat = labels.view(-1)

    # RankingLoss: triplet semi-hard
    ranking_loss = ranking_loss_fnx(visual_emb_flat, text_emb_flat, labels_flat)

    # IdentityLoss: cross-entropy simétrica
    identity_loss = identity_loss_fnx(visual_embeddings, textual_embeddings, labels_flat)

    # Pérdida total: combinación ponderada
    total_loss = (config.ranking_loss_alpha * ranking_loss
                  + config.identity_loss_beta * identity_loss)
\end{minted}

% ============================================================
% Capítulo 10 — Entrenamiento: conceptos + código
% ============================================================
\chapter{Entrenamiento: optimizador, scheduler y mixed precision}
\label{cap:entrenamiento}

Este capítulo explica los conceptos detrás del loop de
entrenamiento y muestra el código línea por línea.

\section{¿Qué es AdamW?}

\subsection{Adam básico}

Adam (\emph{Adaptive Moment Estimation})~\cite{kingma2014adam}
es un optimizador que mantiene dos estadísticas por parámetro:

\begin{itemize}
  \item $m_t$: promedio exponencial del gradiente (momento 1).
  \item $v_t$: promedio exponencial del gradiente al cuadrado
        (momento 2).
\end{itemize}

\begin{align}
  m_t &= \beta_1 m_{t-1} + (1 - \beta_1) g_t \\
  v_t &= \beta_2 v_{t-1} + (1 - \beta_2) g_t^2 \\
  \hat{m}_t &= \frac{m_t}{1 - \beta_1^t} \\
  \hat{v}_t &= \frac{v_t}{1 - \beta_2^t} \\
  \theta_{t+1} &= \theta_t - \eta \cdot \frac{\hat{m}_t}{\sqrt{\hat{v}_t} + \epsilon}
\end{align}

\noindent\textbf{Ventaja}: adapta el learning rate por
parámetro, lo que acelera la convergencia en problemas con
gradientes de magnitudes muy distintas.

\subsection{AdamW}

AdamW~\cite{loshchilov2019adamw} separa el \emph{weight decay}
del optimizador, en vez de aplicarlo como una regularización L2
en la loss. Esto da una regularización más efectiva.

En este proyecto:
\begin{itemize}
  \item $\beta_1 = 0.90$ (\texttt{config.adam\_alpha}).
  \item $\beta_2 = 0.999$ (\texttt{config.adam\_beta}).
  \item lr = 0.001 (\texttt{config.lr}).
\end{itemize}

\section{¿Qué es un LR scheduler?}

Un \emph{learning rate scheduler} modifica el learning rate
durante el entrenamiento. Hay varios tipos:

\begin{itemize}
  \item \textbf{Step}: decae en epochs fijos.
  \item \textbf{Exponential}: decae exponencialmente.
  \item \textbf{Cosine}: cae siguiendo un coseno.
  \item \textbf{ReduceOnPlateau}: decae si la loss no mejora.
\end{itemize}

\noindent Este proyecto usa \texttt{MultiStepLR}: decae por un
factor de 0.1 en los epochs especificados en
\texttt{epoch\_decay = [20, 40]}.

\begin{itemize}
  \item Epochs 1--19: lr = 0.001.
  \item Epochs 20--39: lr = 0.0001.
  \item Epochs 40--60: lr = 0.00001.
\end{itemize}

\noindent\textbf{¿Por qué decaer?} Al principio se quiere
aprender rápido. Cerca del óptimo, lr grande hace oscilar;
reducirlo permite converger.

\section{¿Qué es mixed precision (FP16/BF16)?}

\subsection{Precisiones disponibles}

\begin{itemize}
  \item \textbf{FP32} (32 bits): precisión estándar. Ocupa 4
        bytes por número.
  \item \textbf{FP16} (16 bits): mitad de memoria. Acelera en
        GPUs con Tensor Cores (NVIDIA desde Volta).
  \item \textbf{BF16} (16 bits): mismo rango que FP32 pero
        menos precisión. \emph{No necesita scaler}.
\end{itemize}

\subsection{El problema del FP16}

FP16 tiene rango muy limitado ($\pm 65504$). Si un gradiente es
muy pequeño ($< 6 \cdot 10^{-5}$), se redondea a 0 y
\emph{desaparece}.

\subsection{Solución: GradScaler}

\texttt{torch.cuda.amp.GradScaler} multiplica la loss por un
factor (típicamente $2^{16}$) antes del \texttt{backward}, lo
que efectivamente \emph{amplifica} los gradientes. Luego los
des-escala al hacer \texttt{step}:

\begin{minted}[language=Python, basicstyle=\ttfamily\small, frame=single]{python}
scaler = GradScaler()
scaler.scale(loss).backward()      # loss * factor -> gradientes grandes
scaler.step(optimizer)             # des-escala y aplica
scaler.update()                    # ajusta el factor dinámicamente
\end{minted}

\subsection{Autocast}

\texttt{torch.autocast} aplica casting automático de tipos:

\begin{itemize}
  \item Operaciones \emph{seguras} (matmul, conv): en FP16.
  \item Operaciones \emph{sensibles} (log, softmax, loss): en
        FP32.
\end{itemize}

\noindent Esto se aplica con un context manager:

\begin{minted}[language=Python, basicstyle=\ttfamily\small, frame=single]{python}
with torch.autocast(device_type='cuda', dtype=torch.float16):
    output = model(input)
    loss = criterion(output, target)
\end{minted}

\subsection{GPU vs CPU}

En este proyecto:

\begin{minted}[language=Python, basicstyle=\ttfamily\small, frame=single]{python}
# GPU: FP16 con GradScaler
# CPU: BF16 sin scaler (más estable)
precision_dtype = torch.bfloat16 if config.device == 'cpu' \
                  else torch.float16
\end{minted}

\section{¿Qué hace \texttt{optimizer.zero\_grad()}?}

Por defecto PyTorch \emph{acumula} gradientes entre llamadas a
\texttt{backward()}. Hay que \textbf{limpiar} explícitamente los
gradientes al inicio de cada batch:

\begin{minted}[language=Python, basicstyle=\ttfamily\small, frame=single]{python}
optimizer.zero_grad()  # <- SIEMPRE antes de forward
output = model(input)
loss = criterion(output, target)
loss.backward()         # acumula en .grad
optimizer.step()        # aplica los gradientes
\end{minted}

\noindent\textbf{Alternativa}: \texttt{optimizer.zero\_grad(set\_to\_None=True)}
libera memoria al poner los gradientes en \texttt{None} en vez de
0.

\section{¿Cómo se guarda un checkpoint?}

Un checkpoint es un diccionario con todo lo necesario para
reanudar el entrenamiento:

\begin{minted}[language=Python, basicstyle=\ttfamily\small, frame=single]{python}
checkpoint = {
    'epoch': current_epoch,
    'model_state_dict': model.state_dict(),
    'optimizer_state_dict': optimizer.state_dict(),
    'loss': current_loss
}
torch.save(checkpoint, 'checkpoint.pth.tar')
\end{minted}

\noindent\textbf{state\_dict}: diccionario con todos los
parámetros del modelo (o del optimizer). Es lo que se mueve a
GPU/CPU con \texttt{model.load\_state\_dict()}.

\section{Código: \texttt{train.py} main + loop}

\begin{minted}[language=Python, numbers, firstnumber=34, linerange=34-49]{python}
def parse_args():
    parser = argparse.ArgumentParser(description='Train TextReIDNet on CUHK-PEDES')

    # Fuente del dataset: 'huggingface' (rápido) o 'original' (paper)
    parser.add_argument('--dataset_source', type=str, default='huggingface',
                        choices=['huggingface', 'original'])
    parser.add_argument('--epochs', type=int, default=None)
    parser.add_argument('--batch_size', type=int, default=None)
    parser.add_argument('--lr', type=float, default=None)
    parser.add_argument('--seed', type=int, default=None)
    parser.add_argument('--output_dir', type=str, default='./data/checkpoints')

    return parser.parse_args()
\end{minted}

\begin{minted}[language=Python, numbers, firstnumber=52, linerange=52-110]{python}
def main():
    args = parse_args()

    # Construir config desde config.py
    config = sys_configuration(
        dataset_name="CUHK-PEDES",
        dataset_source=args.dataset_source
    )

    # Aplicar overrides de CLI sobre la config
    if args.epochs is not None:
        config['epoch'] = args.epochs
    if args.batch_size is not None:
        config['batch_size'] = args.batch_size
    if args.lr is not None:
        config['lr'] = args.lr
    if args.seed is not None:
        config['seed'] = args.seed
    if args.output_dir is not None:
        config['model_save_path'] = os.path.abspath(args.output_dir)
        os.makedirs(config['model_save_path'], exist_ok=True)

    # Fijar semilla para reproducibilidad
    set_seed(config.seed)

    # Compartir memoria via filesystem (evita errores con muchos workers)
    torch.multiprocessing.set_sharing_strategy('file_system')

    # GradScaler para mixed precision (FP16/BF16)
    scaler = GradScaler()

    # Construir DataLoaders
    train_data_loader, _, _, train_num_classes = \
        build_cuhkpedes_dataloader(config)

    # Construir modelo y mover a device
    model = TextReIDNet(config).to(config.device)

    # IdentityLoss tiene parámetros entrenables (la capa Linear)
    identity_loss_fnx = IdentityLoss(
        config=config, class_num=train_num_classes
    ).to(config.device)

    # RankingLoss NO tiene parámetros entrenables
    ranking_loss_fnx = RankingLoss(config)

    # Optimizer: AdamW
    optimizer = optim.AdamW(
        model.parameters(),
        betas=(config.adam_alpha, config.adam_beta),
        lr=config.lr
    )

    # Scheduler: decae lr en epochs específicos
    scheduler = optim.lr_scheduler.MultiStepLR(optimizer, config.epoch_decay)
\end{minted}

\begin{minted}[language=Python, numbers, firstnumber=113, linerange=113-191]{python}
for current_epoch in range(1, config.epoch + 1):
    train_ranking_loss_list = []
    train_identity_loss_list = []
    train_total_loss_list = []

    model.train()

    with tqdm(train_data_loader, unit='batch') as tepoch:
        tepoch.set_description(f"Epoch {current_epoch}/{config.epoch}")

        for batch in tepoch:
            # Mover datos al device
            preprocessed_images = batch['preprocessed_images'].to(config.device)
            labels = batch['pids'].to(config.device)
            token_ids = batch['token_ids'].to(config.device)
            orig_token_length = batch['orig_token_lengths'].to(config.device)

            # Limpiar gradientes del batch anterior
            optimizer.zero_grad()

            # Mixed precision
            precision_dtype = torch.bfloat16 if config.device == 'cpu' \
                              else torch.float16
            with torch.autocast(device_type=config.device, dtype=precision_dtype):
                # FORWARD
                visual_embeddings, textual_embeddings = model(
                    image=preprocessed_images,
                    text_ids=token_ids,
                    text_length=orig_token_length
                )

                # Aplanar para losses
                visual_emb_flat = visual_embeddings.view(visual_embeddings.size(0), -1)
                text_emb_flat = textual_embeddings.view(textual_embeddings.size(0), -1)
                labels_flat = labels.view(-1)

                # Calcular losses
                ranking_loss = ranking_loss_fnx(visual_emb_flat, text_emb_flat, labels_flat)
                identity_loss = identity_loss_fnx(visual_embeddings, textual_embeddings, labels_flat)
                total_loss = (config.ranking_loss_alpha * ranking_loss
                              + config.identity_loss_beta * identity_loss)

            # BACKWARD con scaler (para FP16)
            scaler.scale(total_loss).backward()

            # STEP: actualizar parámetros
            scaler.step(optimizer)
            scaler.update()

            # Acumular losses para reporting
            train_ranking_loss_list.append(ranking_loss.item())
            train_identity_loss_list.append(identity_loss.item())
            train_total_loss_list.append(total_loss.item())

            tepoch.set_postfix({
                'rank': f"{np.mean(train_ranking_loss_list):.3f}",
                'id':   f"{np.mean(train_identity_loss_list):.3f}",
                'tot':  f"{np.mean(train_total_loss_list):.3f}",
            })

    scheduler.step()

    # Guardar checkpoint
    ckpt_path = os.path.join(
        config.model_save_path,
        f"TextReIDNet_epoch{current_epoch}.pth.tar"
    )
    torch.save({
        'epoch': current_epoch,
        'model_state_dict': model.state_dict(),
        'optimizer_state_dict': optimizer.state_dict(),
        'loss': np.mean(train_total_loss_list),
    }, ckpt_path)

    latest_path = os.path.join(
        config.model_save_path, "TextReIDNet_latest.pth.tar"
    )
    torch.save({
        'epoch': current_epoch,
        'model_state_dict': model.state_dict(),
        'optimizer_state_dict': optimizer.state_dict(),
        'loss': np.mean(train_total_loss_list),
    }, latest_path)
\end{minted}

\section{Estimación de tiempo}

Con los parámetros default (60 epochs, batch 8, FP16):

\begin{itemize}
  \item \textbf{RTX 2060}: 2--4 días.
  \item \textbf{RTX 3090}: 16--24 horas.
  \item \textbf{A100}: 8--12 horas.
\end{itemize}

\noindent\textbf{Nota}: el HF dataset tiene 179k muestras (vs 40k
del original). El entrenamiento es proporcionalmente más largo
que el paper.

\section{Caveats}

\begin{enumerate}
  \item \textbf{Sin early stopping}: el código entrena los 60
        epochs completos. \texttt{patience = 3} está en config
        pero no se usa.
  \item \textbf{Sin validación intermedia}: no evalúa sobre
        val/test durante el entrenamiento.
  \item \textbf{Sin accumulation de gradientes}: si OOM, hay
        que bajar \texttt{batch\_size}.
  \item \textbf{Checkpoint por epoch}: 61 archivos
        \texttt{.pth.tar} ($\sim$250 MB cada uno para 25M
        parámetros).
\end{enumerate}

% ============================================================
% Capítulo 11 — Evaluación: conceptos + código
% ============================================================
\chapter{Evaluación: extracción de features y métricas}
\label{cap:evaluacion}

Este capítulo explica los conceptos de inferencia y muestra el
código de evaluación. Para la explicación detallada de las
métricas (Top-k, mAP, AP), ver el \cref{cap:metricas}.

\section{¿Qué es la inferencia (evaluación)?}

Una vez entrenado el modelo, queremos medir su rendimiento. El
proceso tiene tres pasos:

\begin{enumerate}
  \item \textbf{Extraer embeddings} de todas las imágenes y
        textos del split test.
  \item \textbf{Computar la matriz de similitud coseno} entre
        imágenes y textos.
  \item \textbf{Calcular métricas} (Top-k, mAP) usando los pids.
\end{enumerate}

\noindent\textbf{Diferencia con entrenamiento}:
\begin{itemize}
  \item En entrenamiento se hace forward + backward.
  \item En inferencia sólo forward (sin gradientes).
\end{itemize}

\section{¿Por qué \texttt{@torch.no\_grad()}?}

\texttt{@torch.no\_grad()} es un decorator que desactiva el
cálculo de gradientes. Beneficios:

\begin{itemize}
  \item \textbf{Menor memoria}: no se almacenan gradientes.
  \item \textbf{Más rápido}: ahorra cómputo.
  \item \textbf{Previene errores}: garantiza que no se
        modifiquen parámetros accidentalmente.
\end{itemize}

\section{¿Qué es \texttt{model.eval()}?}

\texttt{model.eval()} cambia el modo del modelo:

\begin{itemize}
  \item Desactiva \texttt{Dropout} (no se eliminan activaciones).
  \item Usa estadísticas \emph{acumuladas} en \texttt{BatchNorm}
        en lugar de las del batch.
  \item Afecta a capas como \texttt{Dropout}, \texttt{BatchNorm}.
\end{itemize}

\noindent Para inferir SIEMPRE se debe llamar
\texttt{model.eval()} antes.

\section{Galería vs query: convención de este proyecto}

En este proyecto:

\begin{itemize}
  \item \textbf{Query} = texto (descripción de la persona).
  \item \textbf{Gallery} = imágenes.
\end{itemize}

\noindent Para cada query se buscan imágenes similares en la
gallery. La métrica \textbf{Top-k} indica cuántas queries
tienen la imagen correcta en el top-k.

\section{Código: extracción de features de imágenes}

\begin{minted}[language=Python, numbers, firstnumber=42, linerange=42-62]{python}
@torch.no_grad()
def extract_image_features(model, img_loader, device):
    """Extrae embeddings visuales para todas las imágenes del loader."""

    # Modo evaluación (desactiva dropout, etc.)
    model.eval()

    all_features = []
    all_pids = []
    all_paths = []

    for batch in tqdm(img_loader, desc='Extracting image features'):
        if isinstance(batch, (list, tuple)):
            imgs, pids, paths = batch[0], batch[1], batch[2]
        else:
            imgs, pids, paths = batch

        # GPU
        imgs = imgs.to(device)

        # Forward solo en la rama visual
        # image_embedding retorna (B, 1024, 1)
        feats = model.image_embedding(imgs)

        # Aplanar y mover a CPU
        feats = feats.view(feats.size(0), -1).cpu()

        all_features.append(feats)
        all_pids.extend(pids.tolist() if torch.is_tensor(pids) else list(pids))
        all_paths.extend(paths if isinstance(paths, list) else [paths])

    # Concatenar todos los batches
    return torch.cat(all_features, 0), all_pids, all_paths
\end{minted}

\section{Código: extracción de features de textos}

\begin{minted}[language=Python, numbers, firstnumber=65, linerange=65-85]{python}
@torch.no_grad()
def extract_text_features(model, txt_loader, device):
    """Extrae embeddings textuales."""

    model.eval()

    all_features = []
    all_pids = []

    for batch in tqdm(txt_loader, desc='Extracting text features'):
        if isinstance(batch, (list, tuple)):
            pids, token_ids, token_lengths = batch[0], batch[1], batch[2]
        else:
            pids, token_ids, token_lengths = batch

        token_ids = token_ids.to(device)
        token_lengths = token_lengths.to(device)

        # Forward solo en la rama textual
        feats = model.text_embedding(token_ids, token_lengths)
        feats = feats.view(feats.size(0), -1).cpu()

        all_features.append(feats)
        all_pids.extend(pids.tolist() if torch.is_tensor(pids) else list(pids))

    return torch.cat(all_features, 0), all_pids
\end{minted}

\section{Código: similitud coseno}

\begin{minted}[language=Python, numbers, firstnumber=9, linerange=9-17]{python}
def calculate_similarity(image_feature_local, text_feature_local):
    """image: (N_img, 1024), text: (N_txt, 1024)
       returns: (N_img, N_txt) numpy array de similitudes coseno"""

    # Aplanar por si tienen más dims
    image_feature_local = image_feature_local.view(image_feature_local.size(0), -1)
    text_feature_local = text_feature_local.view(text_feature_local.size(0), -1)

    # Normalizar L2 (cada vector a norma 1)
    image_feature_local = image_feature_local / \
        image_feature_local.norm(dim=1, keepdim=True)
    text_feature_local = text_feature_local / \
        text_feature_local.norm(dim=1, keepdim=True)

    # Producto matricial: S[i,j] = cos(img_i, txt_j)
    similarity = torch.mm(image_feature_local, text_feature_local.t())

    # Devolver como numpy (requerido por calculate_ap)
    return similarity.cpu().numpy()
\end{minted}

\section{Código: pipeline completo de evaluación}

\begin{minted}[language=Python, numbers, firstnumber=89, linerange=89-153]{python}
def main():
    args = parse_args()

    config = sys_configuration(
        dataset_name="CUHK-PEDES",
        dataset_source=args.dataset_source
    )
    config['batch_size'] = args.batch_size
    config['device'] = args.device
    config['model_testing_data_split'] = args.split  # 'test' o 'val'

    # Construir dataloaders (devuelve img_loader y txt_loader)
    _, img_loader, txt_loader, _ = build_cuhkpedes_dataloader(config)

    # Construir modelo y cargar checkpoint
    model = TextReIDNet(config).to(config.device)
    ckpt = torch.load(args.checkpoint, map_location=config.device)
    if 'model_state_dict' in ckpt:
        model.load_state_dict(ckpt['model_state_dict'])
    else:
        model.load_state_dict(ckpt)  # compatibilidad

    # PASO 1: extraer features de imágenes
    img_feats, img_pids, img_paths = extract_image_features(
        model, img_loader, config.device
    )

    # PASO 2: extraer features de textos
    txt_feats, txt_pids = extract_text_features(
        model, txt_loader, config.device
    )

    # PASO 3: matriz de similitud coseno
    similarity = calculate_similarity(img_feats, txt_feats)

    # PASO 4: calcular CMC y mAP
    img_pids_t = torch.tensor(img_pids)
    txt_pids_t = torch.tensor(txt_pids)
    cmc, mAP = evaluate(similarity, txt_pids_t, img_pids_t)

    # PASO 5: imprimir resultados
    top1  = cmc[0] * 100
    top5  = cmc[4] * 100 if len(cmc) > 4 else 0
    top10 = cmc[9] * 100 if len(cmc) > 9 else 0

    print(f"  Top-1:  {top1:.2f}%")
    print(f"  Top-5:  {top5:.2f}%")
    print(f"  Top-10: {top10:.2f}%")
    print(f"  mAP:    {mAP*100:.2f}%")

    # Guardar en log
    logger = setup_logger(name='eval_logger',
                           log_file_path=config.test_log_path,
                           write_mode='append')
    logger.info(f"Top-1: {top1:.2f}% | Top-5: {top5:.2f}% | "
                f"Top-10: {top10:.2f}% | mAP: {mAP*100:.2f}%")
\end{minted}

\section{Uso desde la terminal}

\begin{minted}[language=bash, basicstyle=\ttfamily\small, frame=single]{python}
# Evaluar en el split test
.venv/bin/python evaluate.py \
    --checkpoint data/checkpoints/TextReIDNet_latest.pth.tar \
    --dataset_source huggingface \
    --split test

# Cambiar batch size
.venv/bin/python evaluate.py \
    --checkpoint data/checkpoints/TextReIDNet_latest.pth.tar \
    --batch_size 64

# Forzar CPU (lento)
.venv/bin/python evaluate.py \
    --checkpoint data/checkpoints/TextReIDNet_latest.pth.tar \
    --device cpu
\end{minted}

\section{Resultados típicos}

\begin{table}[H]
\centering
\caption{Resultados esperados según el README.}
\label{tab:resultados-tipicos-eval}
\begin{tabular}{lcccc}
\toprule
\textbf{Setup} & \textbf{Top-1} & \textbf{Top-5} & \textbf{Top-10} & \textbf{mAP} \\
\midrule
Sin entrenar (random) & \textasciitilde 0.02\% & 0.05\% & 0.10\% & 0.10\% \\
HF (128×128) 50 batches & 0.02\% & 0.08\% & 0.17\% & 0.10\% \\
HF entrenado completo & 45--52\% & \textasciitilde 70\% & \textasciitilde 80\% & \textasciitilde 40\% \\
Original paper & 54.02\% & 75.51\% & 85.28\% & \textasciitilde 50\% \\
\bottomrule
\end{tabular}
\end{table}

\section{Limitaciones}

\begin{enumerate}
  \item \textbf{Sin gallery augmentation}: cada imagen se evalúa
        una sola vez.
  \item \textbf{Sin query expansion}: no se re-rankea usando
        los top-k.
  \item \textbf{Sin análisis por longitud de caption}.
  \item \textbf{Sin visualización} de aciertos/fallos.
\end{enumerate}

% ============================================================
% Capítulo 12 — Métricas: Top-k, AP, mAP, CMC
% ============================================================
\chapter{Métricas: Top-k, AP, mAP y CMC}
\label{cap:metricas}

Este capítulo desglosa línea por línea cada una de las métricas
que usa el proyecto. Si nunca trabajaste con \emph{re-ID} o
recuperación de imágenes, esta es la lectura obligatoria: vas a
entender qué significa cada número del Top-k y cómo se calcula
exactamente.

\section{El problema de evaluación}

El modelo produce, para cada imagen y cada texto, un vector de
1024 dimensiones (un \emph{embedding}). Para evaluar, comparamos
esos vectores con la \textbf{similitud coseno}:

\begin{equation}
  \text{sim}(\mathbf{I}', \mathbf{T}') =
  \frac{\mathbf{I}' \cdot \mathbf{T}'}
       {\|\mathbf{I}'\| \cdot \|\mathbf{T}'\|}
  \label{eq:sim-cos}
\end{equation}

El resultado es una \textbf{matriz de similitud}
$\mathbf{S} \in \mathbb{R}^{N_{\text{img}} \times N_{\text{txt}}}$
donde $S_{ij}$ indica qué tan parecida es la imagen $i$ al
texto $j$.

\subsection*{Galería vs.\ query}

En este proyecto:
\begin{itemize}
  \item \textbf{Query} = descripción textual (texto).
        Son $N_{\text{txt}} = 29\,930$ en el split test.
  \item \textbf{Gallery} = imágenes.
        Son $N_{\text{img}} = 4\,256$ en el split test.
\end{itemize}

Para cada query (texto) queremos encontrar la imagen correcta
en la gallery. El ranking se obtiene ordenando las imágenes por
similitud descendente.

\section{Top-k accuracy (CMC)}

\subsection*{Qué es}

El \emph{Top-k accuracy} (también llamado \textbf{CMC}, Cumulative
Matching Characteristic) responde a la pregunta:

\begin{quote}
\textit{Dada una descripción textual, ¿la imagen correcta está
entre las primeras $k$ del ranking?}
\end{quote}

\noindent\textbf{Ejemplo:} si para el texto ``\emph{mujer con
camiseta negra y pantalones caqui}'' el ranking de imágenes es:
\begin{enumerate}
  \item img\_382 (correcta)
  \item img\_124
  \item img\_991
\end{enumerate}

\noindent entonces Top-1 = 100\% (acertó de entrada), Top-5 = 100\%
(también está en el top-5). Si la imagen correcta hubiera
aparecido en la posición 7, Top-1 = 0\% pero Top-10 = 100\%.

\subsection*{Cómo se calcula}

Para todos los queries a la vez:

\begin{equation}
  \text{Top-}k = \frac{1}{N_q} \sum_{i=1}^{N_q}
                  \mathbb{1}[\text{correcta}_i \leq k]
  \label{eq:topk}
\end{equation}

\noindent donde $\mathbb{1}[\cdot]$ es la función indicadora (1 si
se cumple, 0 si no).

\section{Average Precision (AP)}

\subsection*{Qué es}

El \acronym{AP}{AP} mide la calidad del ranking \emph{completo},
no sólo del top-1. Es el área bajo la curva
\emph{precisión-recall}.

\subsection*{Ejemplo paso a paso}

Supongamos que para una query hay \textbf{3 imágenes correctas}
en la gallery y el modelo las rankea así (de mayor a menor
similitud):

\begin{table}[H]
\centering
\caption{Ejemplo de ranking para una query con 3 positivos.}
\label{tab:ap-ejemplo}
\begin{tabular}{ccccl}
\toprule
\textbf{Pos.} & \textbf{¿Correcta?} & \textbf{Precisión} & \textbf{Recall} & \textbf{Nota} \\
\midrule
1 & \cmark & 1/1 = 1.00 & 1/3 = 0.33 & primer positivo \\
2 & \xmark & 0/2 = 0.00 & 1/3 = 0.33 & falso positivo \\
3 & \cmark & 1/3 = 0.33 & 2/3 = 0.67 & segundo positivo \\
4 & \xmark & 1/4 = 0.25 & 2/3 = 0.67 & falso positivo \\
5 & \cmark & 2/5 = 0.40 & 3/3 = 1.00 & tercer positivo \\
\bottomrule
\end{tabular}
\end{table}

\noindent\textbf{Precisión} = (correctos vistos) / (total vistos).\\[2pt]
\noindent\textbf{Recall} = (correctos vistos) / (total correctos).

El AP es la integral de la curva precisión-recall. Numéricamente,
la regla del trapecio da:

\begin{align}
  \text{AP} &=
    \sum_{i=1}^{N} \Delta\text{recall}_i \cdot
                 \frac{\text{prec}_i + \text{prec}_{i-1}}{2}
\end{align}

\noindent Aplicado a las posiciones donde hay positivos
(1, 3, 5):

\begin{align*}
  \Delta\text{recall}_1 &= 1/3 \quad \text{prec}_1=1, \text{prec}_0=1 \\
  \Delta\text{recall}_3 &= 1/3 \quad \text{prec}_3=1/3, \text{prec}_2=0 \\
  \Delta\text{recall}_5 &= 1/3 \quad \text{prec}_5=2/5, \text{prec}_4=1/4 \\
  \text{AP} &= \frac{1}{3}(1 + \frac{0.5+0}{2}) +
               \frac{1}{3}(\frac{1}{3}+\frac{1}{4}) \cdot \frac{1}{2} \cdot 2 \\
            &= 0.667 + 0.097 + 0.097 \\
            &= 0.86
\end{align*}

\subsection*{Implementación línea por línea}

\begin{minted}[numbers, firstnumber=20, linerange=20-53]{python}
def calculate_ap(similarity, label_query, label_gallery):
    """
    Args.:
        similarity:     vector 1-D con similitud a cada imagen de la gallery
        label_query:    pid (int) del query
        label_gallery:  vector 1-D con pids de la gallery (mismo orden que similarity)
    Returns: (ap, cmc) o (None, None) si no hay positivos
    """

    # Ordenar índices por similitud DESCENDENTE (mejor primero)
    # np.argsort devuelve ascending; [::-1] invierte el orden
    # np es numpy: import numpy as np
    index = np.argsort(similarity)[::-1]

    # good_index: posiciones en la gallery donde el pid == label_query
    # np.argwhere devuelve índices donde se cumple la condición
    good_index = np.argwhere(label_gallery == label_query)

    # cmc: vector de ceros del mismo tamaño que el ranking
    # Lo llenaremos con 1 desde la posición del primer match
    cmc = np.zeros(index.shape)

    # Máscara booleana: True donde index[i] está en good_index
    # np.in1d chequea pertenencia a un conjunto, eficiente
    mask = np.in1d(index, good_index)

    # Posiciones (en el ranking) donde hay match
    # precision_result[i] = posición en el ranking del i-ésimo match
    precision_result = np.argwhere(mask == True)
    precision_result = precision_result.reshape(
        precision_result.shape[0])

    # Si hay al menos un positivo:
    if precision_result.shape[0] != 0:

        # Marcar con 1 desde el primer match en adelante
        # cmc[k] = 1 significa "el match aparece en posición <= k"
        cmc[int(precision_result[0]):] = 1

        # d_recall: cuánto aumenta el recall por cada positivo
        # 1/N donde N = total de positivos en la gallery
        d_recall = 1.0 / len(precision_result)

        # Acumular AP con regla del trapecio
        ap = 0
        for i in range(len(precision_result)):

            # Precisión en la posición del i-ésimo match
            # = (matches hasta aquí) / (rank del match)
            precision = (i + 1) * 1.0 / (precision_result[i] + 1)

            # Precisión del match anterior (para el trapecio)
            if precision_result[i] != 0:
                old_precision = i * 1.0 / precision_result[i]
            else:
                # Caso especial: primer match en posición 0
                old_precision = 1.0

            # Área del trapecio entre (old_prec, precision) con altura d_recall
            ap += d_recall * (old_precision + precision) / 2

        return ap, cmc

    else:
        # No hay positivos para esta query (no debería pasar)
        return None, None
\end{minted}

\subsection*{Detalles a tener en cuenta}

\begin{itemize}
  \item Si la imagen correcta aparece en la posición 1, AP = 1.0
        (perfect ranking).
  \item Si aparece en la última posición, AP es bajo pero no 0
        (porque se premia que aparezca \emph{al final}, no
        totalmente ausente).
  \item Si una query tiene \emph{un solo} positivo en la gallery,
        AP reduce a Top-1 (porque el recall siempre crece igual).
  \item Si tiene varios positivos, AP captura mejor la calidad
        que Top-1.
\end{itemize}

\section{mAP (mean Average Precision)}

\subsection*{Qué es}

\acronym{mAP}{mAP} es simplemente el promedio de AP sobre todos
los queries:

\begin{equation}
  \text{mAP} = \frac{1}{N_q} \sum_{i=1}^{N_q} \text{AP}_i
  \label{eq:map}
\end{equation}

\noindent donde $N_q$ es el número de queries (textos).

\subsection*{Por qué mAP y no sólo Top-1}

\begin{itemize}
  \item \textbf{Top-1} sólo cuenta si el primero es correcto.
        Ignora si los otros positivos están bien rankeados.
  \item \textbf{mAP} considera el ranking \emph{completo}: si una
        query tiene 5 positivos y todos están en los primeros 5,
        AP = 1.0. Si dos están fuera del top-5, AP < 1.0.
\end{itemize}

\noindent\textbf{En este proyecto} cada query (texto) tiene
\emph{un solo} positivo en la gallery (la imagen correspondiente).
En ese caso mAP y Top-1 están correlacionados pero NO son
iguales:
\begin{itemize}
  \item Top-1 cuenta cuántos queries tienen el positivo en
        posición 1.
  \item mAP promedia la calidad del ranking: un query con el
        positivo en posición 2 da AP = 0.5 (en el caso ideal),
        mientras que Top-1 = 0 para ese query.
\end{itemize}

\section{CMC: el vector completo}

\subsection*{Qué es}

El proyecto devuelve \texttt{cmc} que es un vector de tamaño
$N_{\text{gallery}}$. Para cada posición $k$:

\begin{equation}
  \text{cmc}[k] = \frac{1}{N_q}
                  \sum_{i=1}^{N_q} \mathbb{1}[\text{rank}_i \leq k]
  \label{eq:cmc}
\end{equation}

\noindent Es la fracción acumulada de queries cuyo positivo
aparece en posición $\leq k$. \texttt{cmc[0]} = Top-1,
\texttt{cmc[4]} = Top-5, etc.

\subsection*{Implementación}

\begin{minted}[numbers, firstnumber=56, linerange=56-77]{python}
def evaluate(similarity, label_query, label_gallery):
    """
    similarity:     matriz (N_img, N_txt) de similitud coseno
    label_query:    vector (N_txt,) con pids de los textos (queries)
    label_gallery:  vector (N_img,) con pids de las imágenes (gallery)
    Returns: (cmc, mAP)
    """

    # Convertir de torch.Tensor a numpy array
    label_query = label_query.numpy()
    label_gallery = label_gallery.numpy()

    # Vector acumulador de CMC
    # Tamaño = número de imágenes en la gallery
    # Cada posición k almacena cuántos queries acertaron en top-k
    cmc = np.zeros(label_gallery.shape)

    # Acumulador de AP (para luego sacar promedio = mAP)
    ap = 0

    # Para cada query de texto:
    for i in range(len(label_query)):

        # Calcular AP para esta query
        # similarity[:, i] = columna i = similitud de TODAS las imágenes
        #                    con el texto i-ésimo (orden: filas son imágenes)
        ap_i, cmc_i = calculate_ap(
            similarity[:, i],
            label_query[i],
            label_gallery)

        # Acumular cmc (suma de indicadores)
        cmc += cmc_i

        # Acumular AP
        ap += ap_i

    # Normalizar: cmc[k] = fracción de queries que acertaron en top-k
    cmc = cmc / len(label_query)

    # Normalizar: mAP = promedio de AP sobre todos los queries
    map = ap / len(label_query)

    """
    Comentario sobre cmc:
    cmc[i] es un vector [0, 0, ..., 1, 1, ..., 1] donde el primer 1
    está en la posición del primer match correcto. rank-n y todos
    los ranks posteriores suman un acierto. Al sumar todos los
    vectores y dividir por n_query obtenemos el rank-k accuracy.
    """

    return cmc, map
\end{minted}

\section{Interpretación de números}

\subsection*{Resultados típicos para CUHK-PEDES}

\begin{table}[H]
\centering
\caption{Resultados esperados según el README y el paper.}
\label{tab:resultados-tipicos}
\begin{tabular}{lcccc}
\toprule
\textbf{Setup} & \textbf{Top-1} & \textbf{Top-5} & \textbf{Top-10} & \textbf{mAP} \\
\midrule
Sin entrenar (modelo random) & \textasciitilde 0.02\% & 0.05\% & 0.10\% & 0.10\% \\
HF (128×128) 50 batches     & 0.02\% & 0.08\% & 0.17\% & 0.10\% \\
HF entrenado completo       & 45--52\% & \textasciitilde 70\% & \textasciit\sim 80\% & \textasciitilde 40\% \\
Original paper (paper)      & 54.02\% & 75.51\% & 85.28\% & \textasciitilde 50\% \\
\bottomrule
\end{tabular}
\end{table}

\subsection*{Cómo interpretar un resultado ``Top-1 = 50\%''}

Dice que, dado un texto al azar del test set, en el 50\% de los
casos la imagen correcta aparece en la primera posición del
ranking.

\begin{itemize}
  \item Es ``bueno'' para ReID por texto (un humano也很难
        hacerlo mejor).
  \item Es ``malo'' comparado con face recognition ($>$99\%).
  \item El gap entre Top-1 y Top-5 indica cuántas ``casi
        correctas'' hay en el top-5: si Top-5 = 75\% y
        Top-1 = 50\%, el 25\% restante tiene el correcto en
        posiciones 2--5.
\end{itemize}

\subsection*{Cómo interpretar mAP}

\begin{itemize}
  \item mAP es siempre \textbf{menor o igual} a Top-1
        (matemáticamente: $\text{mAP} \leq \text{cmc}[0]$).
  \item Si Top-1 $\approx$ mAP, el modelo es ``perfecto'' en
        sus predicciones (los positivos casi siempre están en
        posición 1).
  \item Si Top-1 $>$ mAP $\cdot 2$, hay muchas predicciones donde
        el correcto está en posiciones 2--5 (modelo ``casi
        bueno'').
\end{itemize}

\section{Ejemplo numérico completo}

Supongamos 3 queries y 5 imágenes en gallery. La matriz de
similitud es:

\begin{verbatim}
       Q1    Q2    Q3
img1 [ 0.9   0.1   0.3 ]
img2 [ 0.5   0.8   0.2 ]
img3 [ 0.2   0.6   0.7 ]
img4 [ 0.7   0.3   0.5 ]
img5 [ 0.4   0.5   0.4 ]
\end{verbatim}

\noindent y los pids (queries vs gallery):
\begin{itemize}
  \item Q1 busca img2 (pid=2)
  \item Q2 busca img3 (pid=3)
  \item Q3 busca img3 (pid=3)
\end{itemize}

\subsection*{Calcular Top-1}

Para cada query ordenamos la columna de similitud:

\begin{itemize}
  \item \textbf{Q1}: img1 (0.9), img4 (0.7), img2 (0.5), img5 (0.4), img3 (0.2).
        Correcto (img2) en posición \textbf{3}. Top-1: NO.
  \item \textbf{Q2}: img2 (0.8), img3 (0.6), img5 (0.5), img4 (0.3), img1 (0.1).
        Correcto (img3) en posición \textbf{2}. Top-1: NO.
  \item \textbf{Q3}: img3 (0.7), img4 (0.5), img1 (0.3), img2 (0.2), img5 (0.4)
        Correcto (img3) en posición \textbf{1}. Top-1: SÍ.
\end{itemize}

\noindent\textbf{Top-1 = 1/3 = 33.3\%}.

\subsection*{Calcular mAP}

\begin{itemize}
  \item \textbf{AP(Q1)}: positivo en posición 3 de 5.
        AP = 1/3 $\cdot$ (1 + (1/3+0)/2) = 1/3 + 1/9 = 4/9 $\approx$ 0.44.
  \item \textbf{AP(Q2)}: positivo en posición 2 de 5.
        AP = 1/1 $\cdot$ (1 + 0)/2 = 0.5.
  \item \textbf{AP(Q3)}: positivo en posición 1 de 5.
        AP = 1.
\end{itemize}

\noindent\textbf{mAP = (0.44 + 0.5 + 1) / 3 $\approx$ 0.65}.

\subsection*{Calcular Top-5}

Todos los positivos están en posiciones $\leq$ 5 (trivialmente),
así que \textbf{Top-5 = 100\%}.

\section{Resumen}

\begin{table}[H]
\centering
\caption{Resumen de las cuatro métricas.}
\label{tab:resumen-metricas}
\begin{tabular}{ll}
\toprule
\textbf{Métrica} & \textbf{Qué captura} \\
\midrule
Top-1 (cmc[0]) & ¿La imagen correcta es la \#1? \\
Top-5 (cmc[4]) & ¿Está en el top 5? \\
Top-10 (cmc[9]) & ¿Está en el top 10? \\
mAP & Calidad global del ranking \\
\bottomrule
\end{tabular}
\end{table}

\noindent\textbf{Consejo para reportar resultados:} siempre
reportar al menos Top-1, Top-5, Top-10 \emph{y} mAP. Top-1 solo
es insuficiente porque no captura la calidad del ranking más
allá del primer lugar.

% ============================================================
% Capítulo 13 — Apéndices
% ============================================================
\chapter{Apéndices}
\label{cap:apendices}

\section{Glosario}

\begin{description}[leftmargin=2.5cm,style=nextline]
  \item[\textbf{BERT}] \emph{Bidirectional Encoder Representations
    from Transformers}. Modelo de lenguaje preentrenado de Google
    (2018). En este proyecto se usa sólo como tokenizer.

  \item[\textbf{BiGRU}] GRU bidireccional. Variante de RNN que
    procesa la secuencia en ambas direcciones y concatena los
    estados ocultos.

  \item[\textbf{CBIR}] \emph{Content-Based Image Retrieval}.
    Recuperación de imágenes por contenido visual en lugar de
    metadatos.

  \item[\textbf{CMC}] \emph{Cumulative Matching Characteristic}.
    Curva que mide la fracción de queries cuyo positivo aparece
    en el top-k. \texttt{cmc[k]} = Top-(k+1) accuracy.

  \item[\textbf{CUHK-PEDES}] \emph{CUHK Person Description Dataset}.
    Benchmark estándar de \reid{} por texto, 40\,206 imágenes con
    descripciones.

  \item[\textbf{DSC}] \emph{Depthwise Separable Convolution}.
    Factorización de una convolución estándar en depthwise
    (por canal) + pointwise ($1 \times 1$). Reduce parámetros y
    FLOPs.

  \item[\textbf{EfficientNet-B0}] Familia de CNNs con compound
    scaling. B0 es la variante más pequeña (\textasciitilde 5,3M
    parámetros), preentrenada en ImageNet.

  \item[\textbf{FPN}] \emph{Feature Pyramid Network}. Estructura
    que combina features de múltiples escalas para detección.

  \item[\textbf{GRU}] \emph{Gated Recurrent Unit}. RNN con dos
    compuertas (reset, update), más simple que LSTM.

  \item[\textbf{mAP}] \emph{mean Average Precision}. Promedio de
    AP sobre todas las queries. Mide calidad del ranking
    completo.

  \item[\textbf{NMS}] \emph{Non-Maximum Suppression}. Técnica para
    eliminar detecciones redundantes.

  \item[\textbf{Person Re-ID}] Tarea de reconocer a la misma
    persona en diferentes cámaras o momentos.

  \item[\textbf{Triplet loss}] Pérdida que enforce que un anchor
    esté más cerca de un positivo que de un negativo, con un
    margen dado.

  \item[\textbf{U-TextReIDNet}] \emph{Unified Text ReID Network}.
    Modelo del paper que une detección humana con \reid{} por
    texto.

  \item[\textbf{Top-k}] Fracción de queries cuyo positivo aparece
    entre los primeros $k$ resultados.
\end{description}

\section{Tabla maestra de parámetros del modelo}

\begin{longtable}{p{4cm}p{3.5cm}p{6cm}}
\toprule
\textbf{Capa} & \textbf{Forma de salida} & \textbf{Parámetros} \\
\midrule
\endhead
\texttt{Embedding(29610, 512)}
  & \dims{(B, 100, 512)}
  & 15,16 M \\
\texttt{Dropout(0.30)}
  & \dims{(B, 100, 512)}
  & 0 \\
\texttt{BiGRU(512, 1024, bidir)}
  & \dims{(B, 100, 2048)}
  & 4,72 M \\
Suma fwd+bwd / 2
  & \dims{(B, 100, 1024)}
  & 0 \\
Permute + unsqueeze
  & \dims{(B, 1024, 100, 1)}
  & 0 \\
\midrule
\texttt{EfficientNet-B0}
  & \dims{(B, 1280, 16, 16)}
  & 5,29 M \\
\texttt{DSC(1280, 1024)}
  & \dims{(B, 1024, 16, 16)}
  & \textasciitilde 14 K \\
\texttt{AdaptiveMaxPool2d(1, 1)}
  & \dims{(B, 1024, 1, 1)}
  & 0 \\
\texttt{DSC(1024, 1024)}
  & \dims{(B, 1024, 1, 1)}
  & \textasciitilde 10 K \\
\midrule
Classifier(1024, num\_classes)
  & \dims{(B, num\_classes)}
  & \textasciitilde 11 K por clase \\
\midrule
\textbf{Total (sin clasificador)}
  & ---
  & \textbf{\textasciitilde 25,2 M} \\
\textbf{Total (con clasificador)}
  & ---
  & \textbf{\textasciitilde 25,2 M + $1024 \cdot N_c$} \\
\bottomrule
\caption{Recuento completo de parámetros. $N_c$ = número de
identidades (clases) en el dataset de entrenamiento. Para
CUHK-PEDES, $N_c = 11\,003$.}
\label{tab:master-params}
\end{longtable}

\section{Tabla maestra de hiperparámetros de entrenamiento}

\begin{longtable}{p{4cm}p{4cm}p{6cm}}
\toprule
\textbf{Hiperparámetro} & \textbf{Valor} & \textbf{Ubicación} \\
\midrule
\endhead
Épocas
  & 60
  & \fline{config.py}{71} \\
Batch size (default)
  & 8
  & \fline{config.py}{123,129} \\
Batch size (Jetson)
  & 4
  & \fline{config.py}{126} \\
Learning rate inicial
  & 1e-3
  & \fline{config.py}{75} \\
Weight decay
  & (decoupled, AdamW)
  & \fline{train.py}{98} \\
$\beta_1$ AdamW
  & 0.90
  & \fline{config.py}{72} \\
$\beta_2$ AdamW
  & 0.999
  & \fline{config.py}{73} \\
LR decay epochs
  & [20, 40]
  & \fline{config.py}{74} \\
LR decay factor
  & 0.1 (default PyTorch)
  & (implícito) \\
Mixed precision (GPU)
  & FP16
  & \fline{train.py}{130} \\
Mixed precision (CPU)
  & BF16
  & \fline{train.py}{130} \\
Ranking loss weight $\alpha$
  & 1.0
  & \fline{config.py}{87} \\
Identity loss weight $\beta$
  & 1.0
  & \fline{config.py}{88} \\
Triplet margin
  & 0.5
  & \fline{config.py}{86} \\
Seed
  & 3407
  & \fline{config.py}{79} \\
\bottomrule
\caption{Hiperparámetros de entrenamiento.}
\label{tab:master-hyperparams}
\end{longtable}

\section{Resultados reportados en el paper}

\begin{longtable}{p{3.5cm}p{2.5cm}p{2.5cm}p{2.5cm}}
\toprule
\textbf{Dataset} & \textbf{Top-1} & \textbf{Top-5} & \textbf{Top-10} \\
\midrule
\endhead
CUHK-PEDES & 54,02\% & 75,51\% & 85,28\% \\
RSTPReid   & 38,45\% & ---    & ---    \\
\bottomrule
\caption{Resultados del paper original.}
\label{tab:paper-results}
\end{longtable}

\noindent\textbf{U-TextReIDNet}: 38,77 M parámetros totales.

\section{Mapa de archivos}

\begin{verbatim}
textreid-train/
|-- README.md
|-- config.py                     # Configuración global
|-- train.py                      # Script de entrenamiento
|-- evaluate.py                   # Script de evaluación
|-- requirements.txt
|
|-- model/
|   |-- textreidnet.py            # Modelo principal
|   |-- visual_network.py         # EfficientNet-B0 wrapper
|   |-- language_network.py       # BERT tokenizer + BiGRU
|   |-- efficientnet_backbone.py  # Variante multi-stage (no usada)
|   |-- model_utils.py            # DepthwiseSeparableConv, etc.
|
|-- datasets/
|   |-- bases.py                  # ImageTextDataset, etc.
|   |-- cuhkpedes.py              # Loader dataset original
|   |-- cuhkpedes_hf.py           # Loader dataset HuggingFace
|   |-- cuhkpedes_dataloader.py   # Constructor de DataLoaders
|   |-- bert_tokenizer.py         # BERT tokenizer
|   |-- simple_tokenizer.py       # CLIP-style BPE tokenizer
|   |-- tiktoken_tokenizer.py     # OpenAI tiktoken
|
|-- evaluation/
|   |-- ranking_loss.py           # Triplet semi-hard
|   |-- identity_loss.py          # Cross-entropy simétrica
|   |-- focal_loss.py             # Sigmoid focal (no usada)
|   |-- evaluations.py            # Top-k, mAP
|
|-- utils/
|   |-- iotools.py                # read_image, read_json
|   |-- miscellaneous_utils.py    # set_seed, pad_tokens, etc.
|
|-- scripts/
|   `-- download_hf_dataset.py    # Descarga y convierte HF
|
|-- data/
|   `-- CUHK-PEDES-HF/            # Dataset descargado
|       |-- raw.parquet
|       |-- reid_raw.json
|       `-- imgs/all/*.jpg
|
|-- docs/
|   `-- paper.md                  # Paper traducido al español
|
|-- notebooks/
|   `-- 01_lab_validate_dataset.ipynb
|
`-- book/                         # (este libro)
    |-- main.tex
    |-- chapters/*.tex
    `-- figures/*.tikz
\end{verbatim}

\section{Cheatsheet de Bash}

\begin{minted}[language=bash, basicstyle=\ttfamily\small, frame=single]{python}
# Instalar dependencias
pip install -r requirements.txt --extra-index-url https://download.pytorch.org/whl/cu117

# Descargar dataset
python scripts/download_hf_dataset.py

# Entrenar
python train.py --dataset_source huggingface --epochs 60

# Evaluar
python evaluate.py \
    --checkpoint data/checkpoints/TextReIDNet_latest.pth.tar \
    --dataset_source huggingface \
    --split test

# Inspeccionar configuración
python -c "from config import sys_configuration; \
    cfg = sys_configuration(); print(cfg)"
\end{minted}

\section{Referencias}

\begin{itemize}
  \item \textbf{Agyeman, R. \& Rinner, B.} (2024).
    \emph{Decentralized Text-Based Person Re-Identification in
    Multicamera Networks}. IEEE Access.
    DOI: 10.1109/ACCESS.2024.3501382.
    Traducción al español en \file{docs/paper.md}.

  \item \textbf{Tan, M. \& Le, Q. V.} (2019).
    \emph{EfficientNet: Rethinking Model Scaling for
    Convolutional Neural Networks}. ICML.
    arXiv:1905.11946.

  \item \textbf{Devlin, J., Chang, M.-W., Lee, K. \& Toutanova,
    K.} (2018). \emph{BERT: Pre-training of Deep Bidirectional
    Transformers for Language Understanding}. NAACL.
    arXiv:1810.04805.

  \item \textbf{Cho, K. et al.} (2014). \emph{Learning Phrase
    Representations using RNN Encoder-Decoder for Statistical
    Machine Translation}. EMNLP.

  \item \textbf{Howard, A. G. et al.} (2017).
    \emph{MobileNets: Efficient Convolutional Neural Networks
    for Mobile Vision Applications}. arXiv:1704.04861.

  \item \textbf{Li, S., Xiao, T., Li, H., Zhou, B., Yue, D. \&
    Wang, X.} (2017). \emph{Person Search with Natural Language
    Description}. ICCV.

  \item \textbf{Loshchilov, I. \& Hutter, F.} (2019).
    \emph{Decoupled Weight Decay Regularization}. ICLR.
    (AdamW).

  \item \textbf{Lin, T.-Y., Goyal, P., Girshick, R., He, K. \&
    Dollár, P.} (2017). \emph{Focal Loss for Dense Object
    Detection}. ICCV. (RetinaNet).
\end{itemize}


% ============================================================
% Bibliografía
% ============================================================
\backmatter
\bibliographystyle{plainnat}
\bibliography{refs}

\end{document}
